% macro conflict, first definition kept: \newcommand{\Fl}{\mathrm{Fl}}
\documentclass[11pt]{article}
\usepackage{fontspec}
\usepackage[french]{babel}
\usepackage{amsmath}
\usepackage{amssymb}
\usepackage{amsthm}
\usepackage{mathrsfs}
\usepackage{tikz-cd}
\usepackage{enumitem}

\newcommand{\printedpage}[1]{\marginpar{\footnotesize p.~#1}}
\newcommand{\msp}[1]{\marginpar{\footnotesize [#1]}}
\newcommand{\Ch}{\underline{\mathrm{Ch}}}
\newcommand{\Hotp}[1]{\mathrm{Hot}_{\bullet\,#1}}
\newcommand{\hotp}[1]{\mathrm{hot}_{\bullet\,#1}}
\newcommand{\Ob}{\mathrm{Ob}\,}
\newcommand{\HomEns}{\mathrm{Hom}_{\mathrm{Ens}}}
\newcommand{\utpi}{\underset{\sim}{\pi}}
\newcommand{\Fib}{\underline{\mathrm{Fib}}}
\newcommand{\Fl}{\underline{\mathrm{Fl}}}
\newcommand{\Hom}{\underline{\mathrm{Hom}}}
\newcommand{\Cat}{\mathrm{Cat}}
\newcommand{\HOT}{\mathrm{HOT}}
\newcommand{\Hot}{\mathrm{Hot}}
\newcommand{\hot}{\mathrm{hot}}
\newcommand{\id}{\mathrm{id}}
\newcommand{\cM}{\mathcal{M}}
\newcommand{\cF}{\mathcal{F}}
\newcommand{\Ho}{\mathrm{Ho}}
\newcommand{\ho}{\mathrm{ho}}
\newcommand{\Wu}{W^{\mathrm{u}}}
\newcommand{\lc}{\mathrm{lc}}
\newcommand{\circled}[1]{\tikz[baseline=(c.base)]\node[draw,circle,inner sep=1pt](c){#1};}
\newcommand{\deff}{\overset{\text{déf}}{=}}
\begin{document}
% ---- from Der-ChXIIc-30-03-03_p001-005.tex
\printedpage{115}
% head: A. Grothendieck : Dérivateurs, Ch. XII, Édition des Universités de Montpellier II et Paris VII
% foot: Version 30 mars 2003

\msp{page 171}

\section*{VI\quad Critère de $W_\infty$-asphéricité pour les $\Omega^n(X,x)$.}

\subsection*{1\quad Cogitations sur $W_n$.}

On va définir, pour $n \in \mathbf{N}$,
\[
W_n \subseteq \mathrm{Fl}(\mathrm{Cat})
\]
par les conditions suivantes (où $f : X \longrightarrow Y$ est une flèche dans Cat).

\noindent\begin{tabular}{@{}l c l p{0.64\textwidth}@{}}
$f \in W_0$ & $\Longleftrightarrow$ & a) & $\pi_0(f)$ bijectif.\\
 & $\Longleftrightarrow$ & b) & Pour tout préfaisceau d'ensembles $F$ \textit{constant} sur $Y$, $H^0(Y,F) \longrightarrow H^0(X,f^*F)$ est bijectif.\\
 & $\Longleftrightarrow$ & c) & Le foncteur $F \longmapsto f^*F$ des préfaisceaux d'ensembles \textit{constants} sur $Y$ vers les faisceaux d'ensembles constants sur $X$ est pleinement fidèle.
\end{tabular}

\texttt{[Les termes de faisceau ou de préfaisceau sont utilisés ici indifféremment.]}

\textbf{NB} Si $F_0$ est un ensemble, $F_{0\,Y}$ le faisceau d'ensembles constant défini par $F$ sur $Y$, on a
\[
\begin{aligned}
H^0(Y,F_{0\,Y}) &\simeq \HomEns(\pi_0(Y),F_0) && \text{(fonctoriel en $Y$)}\\
\mathrm{Hom}_{Y^\wedge}(F_{0\,Y},G_{0\,Y}) &\simeq \HomEns(\pi_0(Y),\HomEns(F_0,G_0)) && \text{(fonctoriel en $Y$),}
\end{aligned}
\]
ce qui implique que a) $\Longrightarrow$ c) $\underset{\substack{\text{en prenant}\\ \mathrm{Hom}_{Y^\wedge}(e_Y,F)}}{\overset{\text{tautologiquement}}{\Longrightarrow}}$ b) $\Longrightarrow$ a).

\noindent\begin{tabular}{@{}l c l p{0.64\textwidth}@{}}
$f \in W_1$ & $\Longleftrightarrow$ & a) & $\pi_0(f)$ bijectif, et pour tout $x \in \Ob X$, $\pi_1(f,x) : \pi_1(X,x) \longrightarrow \pi_1(Y,f(x))$ est bijectif. (NB Il suffit de le vérifier pour \textit{un} $x$ dans chaque composante connexe de $X$.)\\
 & $\Longleftrightarrow$ & a$'$) & $\Pi(X) \xrightarrow{\Pi f} \Pi(Y)$ est une équivalence de catégories.\\
 & $\Longleftrightarrow$ & b) & $H^0(Y,F) \longrightarrow H^0(X,f^*F)$ bijectif pour tout faisceau d'ensembles \textit{localement constant} sur $Y$, et $H^1(Y,F) \longrightarrow H^1(X,f^*F)$ bijectif si de plus $F$ est un faisceau en groupes.\\
 & $\Longleftrightarrow$ & b$'$) & Comme b), mais pour le $H^1$ on se borne à $F$ \textit{constant}, et on exige $H^1(f)$ \textit{surjectif} seulement.\\
 & $\Longleftrightarrow$ & c) & Le foncteur $F \longmapsto f^*F$ de la catégorie des faisceaux d'ensembles localement constants sur $Y$ vers celle des faisceaux localement constants sur $X$ est une équivalence de catégories.
\end{tabular}

\clearpage

\printedpage{116}
% head: A. Grothendieck : Dérivateurs, Ch. XII, Édition des Universités de Montpellier II et Paris VII
% foot: Version 30 mars 2003

\msp{page 172}

\textbf{NB} Le groupe $\pi_1(X,x)$ est défini via la théorie du groupoïde fondamental. On sait qu'il est en bijection canonique avec $\pi_0(\Omega(X,x))$, où $\Omega(X,x) = \Ch(X;x,x)$ (ce qui permet d'avoir une loi de composition monoïdale sur $\Omega(X,x)$, donnant la loi de groupe sur le $\pi_0$), ou au choix $\Omega(X,x) = \Ch_\infty(X;x,x)$ (plus commode dans certaines situations, mais ne donnant pas directement la loi de composition dans $\pi_0(\Omega(X,x))$).

Je vais poser, pour $n \in \mathbf{N}$,
\[
\pi_n(X,x) = \pi_0(\Omega^n(X,x)),
\]
$\Omega^n$ désignant l'itéré $n$-ième du foncteur $(X,x) \longmapsto \Omega(X,x)$ de la catégorie $\mathrm{Cat}_\bullet$ des objets ponctués de Cat dans elle-même. Je préfère ici travailler avec la théorie \underline{Ch} au lieu de $\Ch_\infty$. Tôt ou tard il faudra donner un énoncé de comparaison \dots

\textbf{Théorème 1} (\footnote[112]{Sans doute \textit{faux}. On doit avoir a) $\Longrightarrow$ b), mais b) $\Longrightarrow$ a) \texttt{[est]} sans doute faux déjà pour $n = 2$ (sauf le cas $X$, $Y$ 1-connexes). Cf. énoncé rectifié \texttt{[p. 178]}.}). \textit{Soit $f : X \longrightarrow Y$ dans} Cat\textit{, et soit $n \in \mathbf{N}$, $n \geq 2$, $S$ une partie de} $\Ob X$ \textit{telle que le composé $S \longrightarrow \Ob X \longrightarrow \pi_0(X)$ soit} \texttt{[un]} \textit{épimorphisme. Les conditions suivantes sont équivalentes.}

\begin{itemize}
\item[\textit{a)}] \textit{$\pi_i(f,x)$ isomorphe pour tout $x \in \Ob X$ et tout $i \in [0,n]$.}

\item[\textit{a$'$)}] \textit{Comme a), mais avec $x \in S$.}

\item[\textit{b)}] \textit{On a}
\[
\begin{array}{rcll}
H^0(Y,F) & \xrightarrow{\ \sim\ } & H^0(X,f^*F) & \textit{pour $F$ faisceau d'ensembles}\\
 & & & \textit{localement constant sur $Y$.}\\
H^1(Y,F) & \xrightarrow{\ \sim\ } & H^1(X,f^*F) & \textit{si de plus $F$ faisceau en groupes.}\\
H^i(Y,F) & \xrightarrow{\ \sim\ } & H^i(X,f^*F) & \textit{si de plus $F$ faisceau abélien, $i \in [2,n]$.}
\end{array}
\]
\end{itemize}

\msp{page 173}

\textbf{NB} L'équivalence de a) et a$'$) devrait être triviale, mais pour cela il faut avoir défini les $\pi_n(X,x)$, pour $x$ variable dans $\Ob X$, comme système local sur $X$, chose que je n'ai pas faite encore. Je vais essayer de le faire dans un courant d'air. Tout d'abord, les $\Omega(X,x)$ sont les fibres d'une $W_\omega$-fibration parfaite (\footnote[113]{\textbf{NB} $W_\omega$ est le plus petit localiseur de Cat (satisfaisant W(1, 2 bis, 3)).}), s'insérant dans le carré cartésien
\[
\begin{tikzcd}
\Ch(X) \arrow[d, "{p_X=(\sigma_X,\beta_X)}"'] & \Omega(X,-) \arrow[l] \arrow[d] \\
X \times X & X. \arrow[l, "\mathrm{diag}_X"'] \arrow[ul, dotted, "i_X"']
\end{tikzcd}
\]

\clearpage

\printedpage{117}
% head: A. Grothendieck : Dérivateurs, Ch. XII, Édition des Universités de Montpellier II et Paris VII
% foot: Version 30 mars 2003

On sait alors que l'on a un foncteur
\[
\begin{aligned}
\Pi X &\longrightarrow \mathrm{Hot}_{W_\omega}\\
x &\longmapsto \mathrm{hot}_{W_\omega}(\Omega(X,x)).
\end{aligned}
\tag*{$(*)$}
\]
Or on a, pour $i \geq 1$,
\[
\pi_i(X,x) = \pi_{i-1}(\ \underbrace{\Omega(X,x)}_{\text{objet ponctué}}\ ).
\]
Cela donne le système local cherché des $\pi_i$, si seulement on avait un foncteur un peu plus précis que $(*)$, savoir
\[
\Pi X \longrightarrow \Hotp{W_\omega} \qquad \text{(catégorie $W_\omega$-homotopique \textit{ponctuée})}
\]
(\footnote[114]{Mais il faut aussi vérifier que les $\pi_n(X,x)$ ($n \geq 1$) sont des foncteurs sur $\Hotp{W_\infty}$, i.e. transforment $W_\infty$-équivalences en itou. Cela résulte de W(4a) pour $W_\infty$, qui implique que $(X,x) \longrightarrow \Omega(X,x)$ se factorise par $\Hotp{W_\infty}$ (en tant que foncteur à valeurs dans $\Hotp{W_\infty}$).}). Or la flèche $i_X : X \longrightarrow \Ch X$ relevant $p_X$ donne une section canonique de $\Omega(X,-)$ sur $X$ -
\msp{page 174}
celle justement qui associe à tout $x \in X$ le point marqué $1_x$ de $\Omega(X;x,x)$ (`lacet vide en $x$'). La question est donc celle-ci :

\begin{quote}
On a une $W$-fibration $\Omega \longrightarrow X$, avec une section donné $\sigma$. Cela définit-il un foncteur canonique % [sic] donné -> donnée, accord avec « section »
\[
X \longrightarrow \Hotp{W},
\]
transformant toute flèche en \texttt{[un]} isomorphisme?
\end{quote}

L'idée naturelle est de prendre
\[
x \longmapsto \hotp{W}(\Omega/x)
\]
en ponctuant $\Omega/x$ par le point marqué $\sigma(x)$ de $\Omega_x \hookrightarrow \Omega/x$. L'ennui, c'est que les morphismes de transition $\Omega/x \longrightarrow \Omega/x'$ (pour $x \xrightarrow{u} x'$ dans $X$) ne sont pas compatibles avec les points marqués. Mais on a une flèche (savoir $\sigma(u)$) allant
\msp{page 175}
de l'image $\sigma(x)/x'$ du point marqué de $\Omega/x$ vers le point marqué $\sigma(x')/x'$ de $\Omega/x'$. Ceci doit suffire pour définir la flèche
\[
(\Omega/x,\sigma(x)/x) \longrightarrow (\Omega/x',\sigma(x')/x')
\]
comme flèche dans $\Hotp{W}$, avec transitivité pour une composée $x \xrightarrow{u} x' \xrightarrow{v} x''$ dans $X$. Je vais pour l'instant me borner
\printedpage{118}
% head: A. Grothendieck : Dérivateurs, Ch. XII, Édition des Universités de Montpellier II et Paris VII
% foot: Version 30 mars 2003
\msp{page 176}
à une indication, quitte à faire des vérifications circonstanciées par la suite. Au besoin, j'utiliserai dès maintenant que pour $i \in \mathbf{N}$, $i \geq 1$,
\[
\pi_i(f,x) : \pi_i(X,x) \longrightarrow \pi_i(Y,f(x))
\]
peuvent s'interpréter comme provenant d'un homomorphisme de systèmes locaux
\[
\utpi_i(X,-) \longrightarrow f^*(\utpi_i(Y,-))
\]
sur $X$. La condition a) du théorème 1 s'interprète en disant que cet homomorphisme est un \textit{isomorphisme} pour $i \in [1,n]$, et la condition a$'$) se borne alors à rappeler qu'il suffit de tester qu'il en est ainsi pour \textit{un} point de chaque composante connexe de $X$.

L'énoncé du théorème 1 garde un sens pour $n = 1$, et dans ce cas a été `rappelé' précédemment (et est $\pm$ tautologique). Cela va nous permettre de procéder par récurrence sur $n$, en supposant que le théorème est prouvé pour $n-1$. Cela signifie que l'on peut supposer que les conditions a) \textit{et} b) sont satisfaites par $f$ pour $i \in [1,n-1]$, et se borner à prouver que, sous ces conditions, la condition a) (ou a$'$), pour $i = n$, équivaut à la condition b) pour $i = n$.

\msp{page 177}

Donc il faut prouver simplement ceci :

\textbf{Lemme} (\footnote[115]{Sans doute faux.}). \textit{Supposons que $\pi_i(f,x)$} \texttt{[est un]} \textit{isomorphisme pour $i \in [0,n-1]$ et b) satisfait pour $i \in [0,n-1]$, où $n \in \mathbf{N}$, $n \geq 2$, est fixé. Alors les conditions suivantes sont équivalentes :}

\begin{itemize}
\item[\textit{a)}] $\pi_n(f,x) : \pi_n(X,x) \longrightarrow \pi_n(Y,f(x))$ \texttt{[est un]} \textit{isomorphisme pour tout $x \in \Ob X$.}

\item[\textit{b)}] $H^n(Y,F) \xrightarrow{\sim} H^n(X,f^*F)$ \texttt{[est un]} \textit{isomorphisme pour tout faisceau abélien localement constant sur $Y$.}
\end{itemize}

Notons que $f$ induit une flèche dans $\mathrm{Cat}_\bullet$,
\[
\Omega(f,x) : \Omega(X,x) \longrightarrow \Omega(Y,f(x)),
\]
d'où
\[
\begin{tikzcd}[row sep=small]
\pi_i(\Omega(X,x)) \arrow[r] \arrow[d, dash, "\wr"'] & \pi_i(\Omega(Y,f(x))) \arrow[d, dash, "\wr"'] \\
\pi_{i+1}(X,x) \arrow[r] & \pi_{i+1}(Y,f(x)),
\end{tikzcd}
\]
s'identifiant à $\pi_{i+1}(f,x)$ :
\[
\pi_i(\Omega(f,x)) \simeq \pi_{i+1}(f,x).
\]

\clearpage

\printedpage{119}
% head: A. Grothendieck : Dérivateurs, Ch. XII, Édition des Universités de Montpellier II et Paris VII
% foot: Version 30 mars 2003

L'hypothèse a) du lemme implique que $\pi_i(\Omega(X,x)) \xrightarrow{\sim} \pi_i(\Omega(Y,y))$ pour $i \in [0,n-1]$. Il s'agit ici des $\pi_i$ basés en les points unités. On devra sans doute prouver que, du fait qu'il s'agit d'un morphisme de monoïdes dans Cat (satisfaisant certaines conditions?) ça reste vrai pour les $\pi_i$ en \textit{tout} point $\xi$ de $\Omega(X,n)$. Par l'hypothèse de récurrence, on saura donc que $\Omega(X,x) \longrightarrow \Omega(Y,y)$ % [sic] \Omega(X,n) -> \Omega(X,x), le point de base est x
\msp{page 178}
induit un isomorphisme pour la cohomologie à coefficients localement constants, jusqu'en dimension $n-1$ inclus.

Finalement, je me suis $\pm$ convaincu que l'énoncé du théorème 1 est faux tel quel. Voici un énoncé que j'éspère vrai, et que je vais essayer de prouver dans un esprit `utilitaire', en utilisant tout ce que je sais en homotopie et cohomologie (et pas seulement ce que j'ai développé `from scratch' dans le modélisateur Cat). % [sic] j'éspère -> j'espère, orthographe

\textbf{Théorème 1 rectifié.} \textit{Considérons, pour un entier $n \geq 2$ donné, les conditions suivantes sur une flèche $f : X \longrightarrow Y$ dans} Cat \textit{(ou une application continue d'espaces topologiques, ou un morphisme dans $\mathbf{\Delta}^\wedge$ etc. etc.) :}

\begin{itemize}
\item[$a_n$)] \textit{Si $Z$ est une `fibre homotopique' de $f$ (NB il y en a essentiellement une par composante connexe de $Y$), alors $Z$ est $n$-connexe, i.e.}
\msp{page 179}
\textit{(p. ex) $\pi_i(Z) = 0$ pour $i \in [0,n]$.}

\item[$b_n$)] \textit{$\pi_0(f)$ est bijectif, et $\forall\ x \in X$, $\pi_i(f,x) : \pi_i(X,x) \longrightarrow \pi_i(Y,f(x))$ est} \texttt{[un]} \textit{isomorphisme pour $i \in [1,n]$. (}\textbf{NB} \textit{Il suffit de le dire pour} \textup{un} \textit{point $x$ de chaque composante connexe de $X$.)}

\item[$c_n$)] $H^i(Y,F) \longrightarrow H^i(X,f^*F)$ \texttt{[est un]} \textit{isomorphisme pour $F$ faisceau localement constant sur $Y$ (avec $F$ abélien si $i \geq 2$, $F$ faisceau en groupes si $i = 1$).}
\end{itemize}

\textit{On a alors}
\[
a_n \ \Longrightarrow\ b_n \ \Longrightarrow \text{(\footnotemark[116])}\ c_n \ \Longrightarrow \text{(\footnotemark[117])}\ a_{n-1}
\]
\footnotetext[116]{si $X$, $Y$ 1-connexes, c'est vraiment prouvé, sinon, il y a des vérifications à faire (qui ont l'air soritales \dots).}%
\footnotetext[117]{\textbf{NB} Même si $X$, $Y$ 1-connexes, il n'est pas clair que $\mathrm{c}_n \Longrightarrow \mathrm{b}_n$, car il faut montrer que $\pi_n(Z) \longrightarrow \pi_n(X)$ est nul, et on trouve seulement que le composé avec $\pi_n(X) \longrightarrow H_n(X)$ est nul. Je présume qu'il y a des contre-exemples.}%
\textit{(et je présume que chacune de ces implications est stricte).}

\textsc{Démonstration}. $\mathrm{a}_n \Longrightarrow \mathrm{b}_n$ est conséquence immédiate de la suite exacte d'homotopie pour une Hot-fibration $\bar{X} \xrightarrow{\bar{f}} Y$ qui factorise $X \xrightarrow{f} Y$ en $f = \bar{f} \circ i$, $i \in W_\infty$. En fait, $\mathrm{a}_n$ \textit{équivaut} à $\mathrm{b}_n$ + la surjectivité des $\pi_{n+1}(X,x) \longrightarrow \pi_{n+1}(Y,f(x))$, laquelle bien sûr ne résulte \textit{pas} de $\mathrm{b}_n$, i.e. de la bijectivité des $\pi_i(f,x)$ pour $i \leq n$ (prendre des $K(\pi,n+1)$). Donc l'implication est bien stricte.

% ---- from Der-ChXIIc-30-03-03_p006-010.tex
\printedpage{120}
% head: A. Grothendieck : Dérivateurs, Ch. XII, Édition des Universités de Montpellier II et Paris VII

\msp{page 180}
$\mathrm{b}_n \Longrightarrow \mathrm{c}_n$. On peut supposer que $f$ est une Hot-fibration (et même Hot-parfait si ça nous chante \dots). On sait déjà que pour
\[
H^i(Y,F) \longrightarrow H^i(X,f^*F) \tag*{$(*)$}
\]
tout est O.K. si $i \in \{0,1\}$. Donc on peut supposer $F$ abélien. Si ça nous avance, on peut supposer, par récurrence, (que l’implication $\mathrm{b}_{n-1} \Longrightarrow \mathrm{c}_{n-1}$ est connue, donc) que $(*)$ est \texttt{[un]} isomorphisme pour $i<n$. Reste à le voir pour $i=n$. On est ramené bien sûr, comme $\pi_0(f)$ est bijectif, au cas où $X$, $Y$ sont 0-connexes, donc il y a essentiellement une seule fibre $Z$. La suite exacte d’homotopie nous apprend que
\[
\pi_i(Z,z)=0 \qquad \text{pour } i \le n-1 \ (Z \text{ est } (n-1)\text{-connexe})
\]
et de plus
\[
\pi_n(Z,z) \longrightarrow \pi_n(X,z) \qquad \text{est nul}
\]
(car $\pi_n(X,z) \longrightarrow \pi_n(Y,y)$ \texttt{[est un]} isomorphisme donc \texttt{[un]} monomorphisme), donc
\[
\pi_{n+1}(Y,y) \longrightarrow \pi_n(Z,z) \qquad \text{\texttt{[est un]} épimorphisme.}
\]
Comme $Z$ est $(n-1)$-connexe (donc aussi 1-connexe), on a
\[
H^i(Z,G)=0\ , \qquad \text{si } 0<i<n\ , \tag*{$(*)$}
\]
pour $G$ localement constant (donc constant) sur $Z$.

\msp{page 181}
Écrivons la suite spectrale de Leray pour $f : X \longrightarrow Y$ et $f^*F = F_X$ :
\[
H^*(X,f^*F) \Longleftarrow E_2^{pq} = H^p(Y,H^q(Z,f^*F)),
\]
où ‘$H^q(Z)$’ doit être vu comme un système local sur $Y$. Donc par $(*)$ on a
\[
E_2^{pq}=0 \qquad \text{si } 0<q<n, \qquad
\begin{tikzpicture}[baseline=0pt, x=0.6cm, y=0.6cm]
\draw[->] (0,0) -- (0,4);
\draw[->] (0,0) -- (6,0);
\node[left] at (0,3) {$n$};
\foreach \i in {0,...,4} { \fill (\i,3) circle (2.5pt); \fill (\i,0) circle (2.5pt); }
\node at (1,1.5) {$0$};
\node at (3,1.5) {$0$};
\end{tikzpicture}
\qquad ,
\]
donc on trouve la suite exacte
\[
0 \longrightarrow H^n(Y,F) \longrightarrow H^n(X,F_X) \xrightarrow{\ \alpha\ } H^0(Y,\ \underbrace{H^n(Z,\overbrace{F_Z}^{M_Z})}_{\begin{array}{l}\scriptstyle = \ \mathrm{Hom}(\pi_n(Z,z),M)\\ \scriptstyle \text{par \texttt{[le]} th. de Hurewicz}\end{array}}\ ) \longrightarrow H^{n+1}(Y,F).
\]

\printedpage{121}
% head: A. Grothendieck : Dérivateurs, Ch. XII, Édition des Universités de Montpellier II et Paris VII

Ici je fixe un point $x$ de $X$, $Z$ désigne la fibre de $f$ en $y=f(x)$, et $M$ celle de $F$ en $y$, $M=F_y$. Le groupe $\pi_1(Y)$ y opère, d’ailleurs $\pi \overset{\text{déf}}{=} \pi_1(X,x) \simeq \pi_1(Y,y)$ aussi. Le $H^0(Y,H^n(Z,F_Z))$ peut s’interpréter, sauf erreur, comme
\[
\mathrm{Hom}_\pi(\pi_n(Z,x),M),
\]
en tenant compte de l’opération de $\pi=\pi_1(X,x)$ sur $\pi_n(Z,x)$ également. (À vérifier.) Il faut interpréter la flèche
\[
H^n(X,F_X) \longrightarrow \mathrm{Hom}_\pi(\pi_n(Z,x),M),
\]
via l’application composée \texttt{[où $z=x$, vu comme point de $Z$]}
\[
\pi_n(Z,z) \longrightarrow \pi_n(X,x) \longrightarrow H_n(X,\mathbf{Z}). \tag*{$(*)$}
\]
Si $F$ est \textit{constant} sur $Y$ (p. ex. si $Y$ \texttt{[est]} 1-connexe), donc de valeur $M$, on a les flèches canoniques
\[
H^n(X,F_X) = H^n(X,M) \longrightarrow \mathrm{Hom}(H_n(X,\mathbf{Z}),M) \longrightarrow \mathrm{Hom}(\pi_n(Z,z),M)\text{(\footnotemark[118])} \tag*{$(*')$}
\]
\footnotetext[118]{grâce à la composition avec $(*)$.}
\texttt{[lire : }$\mathrm{Hom}_\pi$\texttt{]},

\msp{page 182}
et sûrement (en dehors de toute hypothèse de type $\mathrm{b}_n$ ou autre) l’application canonique directe ‘de localisation sur $Y$’, suivi de \dots,
\[
H^n(X,F_X) \longrightarrow H^0(Y,H^n(Z,F_Z)) \longrightarrow \mathrm{Hom}_\pi(\pi_n(Z),M) \tag*{$(**)$}
\]
n’est autre que la composée des applications canoniques ci-dessus. Dans le cas actuel où on suppose $\mathrm{b}_n$, la deuxième flèche dans $(**)$ est \texttt{[un]} isomorphisme, donc la composée ‘n’est autre’ que la flèche de localisation elle-même. D’autre part $\pi_n(Z,z) \longrightarrow \pi_n(X,x)$ dans $(*)$ est nul, donc aussi $\pi_n(Z,z) \longrightarrow H_n(X,\mathbf{Z})$, donc aussi la flèche composée $(*')$ qui s’en déduit. Revenons alors à la suite exacte de Leray (page précédente), on trouve que $H^n(Y,F) \xrightarrow{\ \sim\ } H^n(X,F_X)$.

Comment récupérer le cas où $F$ n’est \textit{pas} constant? Dans le cas $n=2$, cela marchait, en passant aux revêtements universels $\tilde{X}$, $\tilde{Y}$ de $X$, $Y$ (par rapport aux points $x$, $y$ choisis), et en écrivant les suites spectrales
\[
\begin{tikzcd}[row sep=small, column sep=small]
H^*(Y,F) \arrow[d] & \Longleftarrow & E_2^{pq} \arrow[d] & = & H^p(\pi,H^q(\tilde{Y},F_{\tilde{Y}})) \\
H^*(X,F_X) & \Longleftarrow & E^{\prime pq}_2 & = & H^p(\pi,H^q(\tilde{X},F_{\tilde{X}})).
\end{tikzcd}
\]
L’hypothèse $\mathrm{b}_n$ est conservée en passant de $f : X \longrightarrow Y$ à $\tilde{f} = \tilde{f}_x : \tilde{X} \longrightarrow \tilde{Y}$, donc on sait que
\printedpage{122}
% head: A. Grothendieck : Dérivateurs, Ch. XII, Édition des Universités de Montpellier II et Paris VII
\msp{page 183}
\[
E_2^{pq} \xrightarrow{\ \sim\ } E^{\prime pq}_2 \qquad \text{\texttt{[est un]} isomorphisme si } q \le n. \tag*{$(*)$}
\]
Dans le cas $n=2$, utilisant que
\[
E_2^{p1} = E^{\prime p1}_2 = 0 \qquad \text{car } H^1(\tilde{X}) = H^1(\tilde{Y}) = 0, \qquad
\begin{tikzpicture}[baseline=0pt, x=0.5cm, y=0.5cm]
\draw[->] (0,0) -- (0,2.8);
\draw[->] (0,0) -- (5,0);
\node[above left] at (0,2.6) {$q$};
\node[left] at (0,2) {$2$};
\foreach \i in {0,...,3} { \fill (\i,2) circle (2.5pt); \fill (\i,0) circle (2.5pt); }
\node at (0,1) {$0$};
\node at (1,1) {$0$};
\node at (2,1) {$0$};
\end{tikzpicture}
\qquad ,
\]
on trouve deux suites exactes
\[
\begin{tikzcd}
0 \arrow[r] & H^2(\pi,H^0(\tilde{Y})) \arrow[r] \arrow[d, "\text{iso}"] & H^2(Y) \arrow[r] \arrow[d] & H^0(\pi,H^2(\tilde{Y})) \arrow[r] \arrow[d, "\text{iso}"] & H^3(\pi,H^0(\tilde{Y})) \arrow[d, "\text{iso}"] \\
0 \arrow[r] & H^2(\pi,H^0(\tilde{X})) \arrow[r] & H^2(X) \arrow[r] & H^0(\pi,H^2(\tilde{X})) \arrow[r] & H^3(\pi,H^0(\tilde{X})),
\end{tikzcd}
\]
et un homomorphisme de l’une dans l’autre, qui permet de conclure par le lemme des 5 que $H^2(Y,F) \longrightarrow H^2(X,F_X)$ est \texttt{[un]} isomorphisme. Mais pour $n$ quelconque, la relation $(*)$ sur les termes $E_2$ (pour $q \le n$) ne me semble pas impliquer la relation similaire pour les $E_n$ supérieures, et de là pour les termes $E_\infty$, ce qui permettrait de conclure.

Je reviens donc à la suite exacte p. 181, avec $F$ faisceau \textit{tordu.} J’ai envie de trouver une application canonique
\[
H^n(X,G) \longrightarrow \overbrace{\mathrm{Hom}_\pi(\pi_n(X,x),G_x)}^{\text{\scriptsize Homs \textit{additifs} bien sûr, i.e. de $\mathbf{Z}[\pi]$-modules}} \qquad ,
\]
où $G$ est un faisceau abélien localement constant quelconque sur $X$, et $\pi = \pi_1(X,x)$. Si j’interprète
\msp{page 184}
$\pi_n(X,x)$ par des classes d’homotopie d’applications ponctuées
\[
(S^n,x_0) \xrightarrow{\ \varphi\ } (X,x)
\]
j’ai bien, pour une telle application $\varphi$,
\[
\varphi^* : H^n(X,G) \longrightarrow H^n(S^n,\ \overbrace{\varphi^*G}^{\simeq\,(G_x)_{S^n}}\ ) \simeq G_x
\]
(tenant compte que $S^n$ est 1-connexe, et que sa $n$-cohomologie à valeurs dans un faisceau constant $M$ est $M$.) Sûrement l’application
\[
H^n(X,G) \longrightarrow G_x
\]
\printedpage{123}
% head: A. Grothendieck : Dérivateurs, Ch. XII, Édition des Universités de Montpellier II et Paris VII
ainsi définie ne dépend que de la classe d’homotopie ponctuée de $\varphi$, d’où
\[
\pi_n(X,x) \times H^n(X,G) \longrightarrow G_x,
\]
et il faut de plus prouver a) que c’est compatible avec les opérations de $\pi = \pi_1(X,x)$, et b) que c’est additif en $\pi_n(X,x)$, pour avoir l’application cherchée
\[
H^n(X,G) \longrightarrow \mathrm{Hom}_\pi(\pi_n(X,x),G_x).
\]
Si j’admets ceci, il n’y a qu’un pas pour admettre aussi que la flèche qui nous intéresse vraiment (dans la suite exacte p. 181)
\msp{page 185}
\[
H^n(X,F_X) \longrightarrow \mathrm{Hom}_\pi(\pi_n(Z,z),\underbrace{M}_{=\,F_y})
\]
n’est autre que la composée de la flèche précédente (page précédente) pour $G = F_X$, avec
\[
\mathrm{Hom}_\pi(\pi_n(X,x),M) \longrightarrow \mathrm{Hom}_\pi(\pi_n(Z,x),M)
\]
déduite de
\[
\pi_n(Z,x) \longrightarrow \pi_n(X,x).
\]
Or cette flèche est nulle, par l’hypothèse $\mathrm{b}_n$. Donc la flèche $\alpha$ dans la suite exacte page 181 est nulle, donc $H^n(Y,F) \longrightarrow H^n(X,F_X)$ est \texttt{[un]} isomorphisme, q.e.d. (mais il y a pas mal de travail de vérification à faire \dots).

$\mathrm{c}_n \Longrightarrow \mathrm{a}_{n-1}$. On peut procéder par récurrence sur $n$. (NB L’enoncé a un sens même pour $n=1$, où il est prouvé, et même $\mathrm{c}_1 \Longrightarrow \mathrm{b}_1$, et $\mathrm{b}_1 \Longrightarrow \mathrm{a}_0$, plus généralement $\mathrm{b}_n \Longrightarrow \mathrm{a}_{n-1}$, immédiat par \texttt{[la]} suite exacte d’homotopie.) Donc on sait déjà, en plus de l’hypothèse $\mathrm{c}_n$, qu’on a $\mathrm{a}_{n-2}$, i.e. $Z$ est $(n-2)$-connexe. Il faut prouver qu’il est même $(n-1)$-connexe, i.e. $\pi_{n-1}(Z)=0$. Le cas $n=2$ demande % [sic] L’enoncé -> L’énoncé, missing accent on é
\msp{page 186}
peut-être une attention particulière, car on ne sait pas même encore que $Z$ est 1-connexe. Mais on sait que $\mathrm{c}_2$ implique $\mathrm{c}_1$, qui équivaut \texttt{[?]} à $\mathrm{b}_1$, donc on a
\[
\pi_1(X,x) \longrightarrow \pi_1(Y,y) \qquad \text{isomorphisme,}
\]
donc
\[
\pi_1(Z,x) \longrightarrow \pi_1(X,x) \qquad \text{\texttt{[est]} \textit{nul,}}
\]
donc
\[
\pi_2(Y,y) \longrightarrow \pi_1(Z,x) \qquad \text{\texttt{[est un]} épimorphisme,}
\]
ce qui implique tout au moins que $\pi_1(Z,x)$ est \textit{abélien.} Il faut le voir comme un $\pi$-module, et prouver sa nullité comme tel (moyennant $\mathrm{c}_2$).

\printedpage{124}
% head: A. Grothendieck : Dérivateurs, Ch. XII, Édition des Universités de Montpellier II et Paris VII

Passons au cas général. On a pour $F$ localement constant abélien sur $Y$
\[
H^*(X,F_X) \Longleftarrow \underbrace{E_2^{pq} = H^p(Y,H^q(Z,F_Z))}_{\begin{array}{l}\scriptstyle = 0 \text{ si } 0<q<n-1\\ \scriptstyle \text{car } Z \text{ est } (n-2)\text{-connexe}\end{array}}, \qquad
\begin{tikzpicture}[baseline=0pt, x=0.6cm, y=0.6cm]
\draw[->] (0,0) -- (0,5);
\draw[->] (0,0) -- (6,0);
\node[left] at (0,4) {$n-1$};
\foreach \i in {0,...,4} { \fill (\i,4) circle (2.5pt); \fill (\i,0) circle (2.5pt); }
\node at (1,2) {$0$};
\node at (3,2) {$0$};
\end{tikzpicture}
\qquad ,
\]
D’où
\[
0 \longrightarrow H^{n-1}(Y,F) \longrightarrow H^{n-1}(X,F_X) \longrightarrow \underbrace{H^0(Y,H^{n-1}(Z,F_Z))}_{\begin{array}{l}\scriptstyle \mathrm{Hom}_\pi(\pi_{n-1}(Z,x),F_x)\\ \scriptstyle \text{par Hurewicz}\end{array}} \longrightarrow H^n(Y,F) \longrightarrow H^n(X,F_X).
\]
L’hypothèse $\mathrm{c}_n$ implique qu’on a
\msp{page 187}
\[
\mathrm{Hom}_\pi(\pi_{n-1}(Z,x),\underbrace{M}_{=\,F_x}) = 0
\]
et ceci \textit{quel que soit le $\pi$-module $M$,} puisque $\pi = \pi_1(X,x) \xrightarrow{\ \sim\ } \pi_1(Y,y)$. Cela implique (en prenant p. ex. $M = \pi_{n-1}(Z,x)$ !) que $\pi_{n-1}(Z,x)=0$. (Mais en faisant attention, dans le cas $n=2$, qu’on sait déjà que $\pi_1(Z,x)$ est abélien).

Donc le théorème rectifie page 178 a bien l’air vrai. Il est un peu triste, puisque dans chacune des implications % [sic] rectifie -> rectifié, past participle intended
\[
\mathrm{a}_n \ \Longrightarrow\ \mathrm{b}_n \ \Longrightarrow\ \mathrm{c}_n \ \Longrightarrow\ \mathrm{a}_{n-1}
\]
on perd un chouia, et en fin de compte on sort avec la magnifique implication
\[
\mathrm{a}_n \ \Longrightarrow\ \mathrm{a}_{n-1}.
\]
De toutes façons, il est $\pm$ clair que le “bon” localiseur $W_n$, celui qui doit définir les “$n$-types d’homotopie” comme objets de $\mathrm{Hot}_{W_n}$, est défini par la condition $\mathrm{a}_n$.

\section*{2\quad Critère de $W_n$-asphéricité.}

Je vais essayer de faire mieux quand $Y$ est réduit à un point. Cette fois, je

% ---- from Der-ChXIIc-30-03-03_p011-015.tex
\printedpage{125}
% head: A. Grothendieck : Dérivateurs, Ch. XII, Édition des Universités de Montpellier II et Paris VII
% foot: Version 30 mars 2003

\msp{page 188}
veux être strict, et donner (si je peux) des démonstrations complètes dans le cadre où je me place à présent, celui du modéliseur Cat, et en n'utilisant que ce que je sais y démontrer directement à présent, sans me référer à la théorie semi-simpliciale ou topologique.

\textbf{Théorème 2} (?). \textit{Soit $X$ dans} Cat, \textit{$n \in \mathbf{N}$ un entier $\geq 1$. Alors les conditions suivantes sont équivalentes :}

\begin{itemize}
\item[\textit{$a_n$)}] \textit{$X$ est $0$-connexe, et si $x$ est dans $X$, $\pi_i(X,x) = 0$ pour $i \in [1,n]$ (condition qui ne dépend pas du choix de $x$), i.e. les $\Omega^i(X,x)$ pour $i \in [1,n]$ sont $0$-connexes.}
\item[\textit{$b_n$)}] \textit{Pour tout groupe $G$, $G \longrightarrow H^0(X,G_X)$ est un isomorphisme, et $H^1(X,G_X) = 0$. De plus, si $G$ est abélien, on a aussi $H^i(X,G_X) = 0$, pour $i \in [1,n]$.}
\end{itemize}

\textsc{Démonstration}. Pour $n = 1$ cette équivalence est $\pm$ tautologique et a été rappelée (sous une forme légèrement différente) au début (p.~71). On procède par récurrence. On peut donc supposer que $n \geq 2$, que $X$ satisfait $\mathrm{a}_{n-1}$, $\mathrm{b}_{n-1}$, et prouver qu'alors $\mathrm{a}_n$, $\mathrm{b}_n$ sont équivalentes, i.e. $\pi_n(X,x) = 0$ équivaut à $H^n(X,M) = 0$ pour tout groupe abélien $M$.

\msp{page 189}
On utilise la $W_\infty$-fibration (parfaite)
\[
\begin{tikzcd}[row sep=large]
E = \Ch(X;x,-) \arrow[d] & E\ W_\infty\text{-asphérique} \\
X &
\end{tikzcd}
\]
de fibre $Z$ en $x$ égale à $\Omega(X,x)$.

Supposons $\mathrm{a}_n$) pour $X$, donc $a_{n-1}$) pour $\Omega(X,x)$. Par hypothèse de récurrence, on a donc
\[
H^i(Z,M) = 0 \qquad \text{pour } i \in [1,n-1],
\]
$M$ groupe abélien. Donc dans la suite spectrale de Leray,
\[
H^*(E,M) \Longleftarrow E_2^{pq} = H^p(X,H^q(Z,M))
\]
on a
\[
E_2^{pq} = 0 \qquad \text{si } q \in [1,n-1], \qquad
\begin{tikzpicture}[baseline=0]
\draw[->] (0,0) -- (2.2,0);
\draw[->] (0,0) -- (0,2.1);
\foreach \x in {0,0.35,0.7,1.05,1.4} {
  \fill (\x,0) circle (1.5pt);
  \fill (\x,1.4) circle (1.5pt);
}
\node[left] at (-0.05,1.4) {$n$};
\node at (0.35,0.7) {$0$};
\node at (1.05,0.7) {$0$};
\end{tikzpicture}
\quad .
\]

\printedpage{126}
% head: A. Grothendieck : Dérivateurs, Ch. XII, Édition des Universités de Montpellier II et Paris VII
% foot: Version 30 mars 2003

On a donc la suite exacte
\[
0 \longrightarrow H^n(X,M) \longrightarrow H^n(E,M)
\]
et comme $E$ est $W_\infty$-asphérique, on a $H^n(E,M) = 0$, donc $H^n(X,M) = 0$.

Inversement, supposons $\mathrm{b}_n$), prouvons $\mathrm{a}_n$). On sait déjà qu'on a $\mathrm{a}_{n-1}$), donc $\mathrm{a}_{n-2}$ pour $\Omega(X,x)$, ce qui implique
\[
H^i(Z,M) = 0 \qquad \text{pour } i \in [1,n-2],
\]
donc
\[
E_2^{pq} = 0 \qquad \text{si } q \in [1,n-2].
\]
On a donc la suite exacte

\msp{page 190}
\[
0 \longrightarrow \underbrace{H^{n-1}(X,M)}_{=\,0} \longrightarrow \underbrace{H^{n-1}(E,M)}_{=\,0} \longrightarrow H^0(X,H^{n-1}(Z,M)) \longrightarrow \underbrace{H^n(X,M)}_{=\,0} \longrightarrow H^n(E,M),
\]
donc
\[
H^0(X,\ \underbrace{H^{n-1}(Z,M)}_{\begin{array}{l}\text{système local}\\ \textit{constant}\text{ sur } X,\\ \text{car } X \text{ est 1-connexe}\end{array}}\ ) = 0,
\]
donc
\[
H^{n-1}(Z,M) = 0,
\]
ce qui, par hypothèse de récurrence, implique $\pi_{n-1}(Z,z) = 0$, donc $\pi_n(X,x) = 0$, q.e.d.

\textbf{Corollaire.} \textit{Pour que $X$ soit $W_\infty$-asphérique, i.e. pour qu'on ait $G \xrightarrow{\ \sim\ } H^0(X,G)$, $H^1(X,G) = 0$} \texttt{[pour]} \textit{tout groupe} \texttt{[}$G$\texttt{]}\textit{, et $H^i(X,G) = 0$ si $i \geq 2$ pour $G$ abélien, il faut et il suffit que} $1^\circ$) \textit{$X$ soit} $0$\textit{-connexe, et ($x$ étant un point de $X$)} $2^\circ$) \textit{les $\Omega^i(X,x)$} \texttt{[soient]} $0$\textit{-connexes pour $i \geq 1$.}

\msp{page 191}

\section*{3\quad Sur la philosophie des fibres homotopiques.}

Le critère d'asphéricité (pour $W_\infty$) que je viens d'obtenir, grâce aux foncteur $\Ch(X)$ (ou au choix, $\Ch_\infty(X)$) et ses dérivés $\Omega(X,x)$ et les itérés, enlève tout doute (s'il pouvait en rester encore) que les `espaces de chemins' que j'ai construit (assez péniblement!) sont bel et bien `les bons'. Je voudrais maintenant me rendre compte de façon la plus précise possible de leur `canonicité', dans quelle mesure ces constructions sont indépendant d'un choix (p.ex. celui entre la théorie $\Ch$, et $\Ch_\infty$). % [sic] aux foncteur -> au foncteur, number agreement

Tout revient, finalement, à des questions de factorisation d'une flèche quelconque
\[
f : X \longrightarrow Y
\]
\printedpage{127}
% head: A. Grothendieck : Dérivateurs, Ch. XII, Édition des Universités de Montpellier II et Paris VII
% foot: Version 30 mars 2003
de Cat en
\[
X \xrightarrow{\ i_f\ } \bar{X} \xrightarrow{\ \bar{f}\ } Y,
\]
avec
\[
i_f \in W, \qquad \bar{f} \in W\text{-fib}
\]
(le fait que $\bar{f}$ puisse même être choisie $W_\omega$-parfait, ne me paraissant pas essentiel ici).

C'est bien de cela qu'il s'agit pour construire

\msp{page 192}
un objet `chemins dans $X$'
\[
\Ch(X),
\]
en factorisant l'application diagonale
\[
X \xrightarrow{\ \mathrm{diag}_X\ } X \times X,
\]
et aussi quand on construit $\Omega(X,x)$ comme fibre en $x$ d'une $W$-fibration
\[
E(X,x) \longrightarrow (X,x)
\]
qui factorise l'inclusion
\[
\{x\} \longrightarrow X
\]
(donc $E(X,x)$ étant $W$-asphérique).

On peut être plus ou moins exigeant, et regarder, soit pour un `but' $Y$ (ou plutôt un $S$) fixé dans Cat, une factorisation fonctorielle
\[
C : \Cat/S \longrightarrow \Fib_W S \hookrightarrow \Cat/S
\]
du morphisme structural $X \longrightarrow S$ d'un objet, en
\[
X \xrightarrow{\ i_X\ } C(X) \longrightarrow S,
\]
i.e. on se donne le foncteur $C$, \textit{plus}

\msp{page 193}
le morphisme fonctoriel
\[
i_C : \mathrm{id}_{\Cat/S} \longrightarrow C
\]
avec
\[
(i_X : X \longrightarrow C(X)) \in W \qquad \forall\ X \text{ dans } \Cat/S,
\]
soit être encore plus exigeant et faire varier en même temps la source et le but de $f : X \longrightarrow Y (= S)$, et demander un foncteur
\[
\begin{aligned}
C : \Fl(\Cat) \quad &\longrightarrow \quad \Cat \\
f \quad &\longmapsto \quad C(f),
\end{aligned}
\]
\printedpage{128}
% head: A. Grothendieck : Dérivateurs, Ch. XII, Édition des Universités de Montpellier II et Paris VII
% foot: Version 30 mars 2003
avec des morphismes fonctoriels en $f$
\[
\underbrace{s(f)}_{=\,X} \xrightarrow{\ i_f\ } \underbrace{C(f)}_{=\,\bar{X}} \xrightarrow{\ p_f\ } \underbrace{b(f)}_{=\,Y}
\]
($i$, $p$ des flèches dans $\Hom(\Fl(\Cat),\Cat)$), avec
\[
i_f \in W, \qquad p_f \in W\text{-fib}
\]
pour toute flèche $f$ dans Cat.

\msp{page 194}
Dans les deux points de vue, dans quel sens a-t-on \textit{unicité} d'une telle construction?

Prenons d'abord le premier point de vue, factorisation des flèches à but fixé $S$, donc \texttt{[c'est une]} question dans $\Cat/S$. J'ai oublié une propriété importante, mais elle résulte de ce qui a été dit : le foncteur
\[
\begin{tikzcd}[column sep=small]
X \longmapsto C(X) : & \Cat/S \arrow[rr] \arrow[dr, dotted] & & \Cat/S \\
& & \Fib_W S \arrow[ur, dotted, hook] &
\end{tikzcd}
\]
\textit{transforme} \texttt{[les]} \textit{flèches} \texttt{[de]} \textit{$W$ en itou}, comme on voit sur le diagramme commutatif
\[
\begin{tikzcd}[row sep=large, column sep=large]
X \arrow[r, "i_X"] \arrow[d, "f"'] & C(X) \arrow[d, "C(f)"] \\
Y \arrow[r, "i_Y"] & C(Y),
\end{tikzcd}
\]
où $i_X$, $i_Y$ sont dans $W$, donc
\[
f \in W \qquad \Longleftrightarrow \qquad C(f) \in W.
\]
Soit
\[
(C',i')
\]
un deuxième foncteur, $\Cat/S \longrightarrow \Cat/S$,

\msp{page 195}
muni d'un morphisme
\[
i' : \mathrm{id}_{\Cat/S} \longrightarrow C'
\]
qui est dans $W$ argument par argument, et tel que $C'$ se factorise par $\Fib_W S$. On peut comparer $C$ et $C'$ moyennant $C \circ C'$, car
\[
\begin{tikzcd}[column sep=small]
& C \circ C' & \\
C = C \circ \mathrm{id} \arrow[ur, "C\ast i'"] & & C' = \mathrm{id} \circ C', \arrow[ul, "i\ast C'"']
\end{tikzcd}
\]
\printedpage{129}
% head: A. Grothendieck : Dérivateurs, Ch. XII, Édition des Universités de Montpellier II et Paris VII
% foot: Version 30 mars 2003
i.e. on a, pour tout $X$ dans $\Cat/S$,
\[
\begin{tikzcd}[column sep=small]
& C(C'(X)) & \\
C(X) \arrow[ur, "C(i'_X)"] & & C'(X), \arrow[ul, "i_{C'(X)}"']
\end{tikzcd}
\]
le deux flèches étant dans $W$ (celle de gauche parce que $C(W) \subseteq W$). Ce sont des flèches dans $\Fib_W S$, et si on suppose que $W$ satisfait W(4b), elles induisent des $W$-équivalences en chaque fibre. Cela montre donc déjà % [sic] le deux flèches -> les deux flèches, typo

\msp{page 196}
qu'à isomorphisme (peut-être non-canonique) près, le foncteur composé défini par $C$,
\[
\begin{tikzcd}
X \arrow[r, mapsto] & C(X) \arrow[r, mapsto] & (s \longmapsto \mathrm{hot}_W(C(X)_s)) \\
\Cat/S \arrow[r] \arrow[rr, bend right, "\varphi_C"'] & \Fib_W S \arrow[r] & \Hom(S^\circ,\mathrm{Hot}_W),
\end{tikzcd}
\]
ne dépend pas du choix particulier de $C$, i.e. deux tels foncteurs `fibres homotopiques' (correspondant à deux choix $C$, $C'$) sont toujours isomorphes. Mais au lieu d'utiliser $C \circ C'$ pour définir un tel isomorphisme, on aurai pu utiliser $C' \circ C$, ou quelque autre procédé. On aimerait comprendre s'il y a unicité véritable, une caractérisation intrinsèque de ces isomorphismes $\varphi_C \xrightarrow{\ \sim\ } \varphi_{C'}$, indépendamment de telle construction ou de telle autre. % [sic] aurai pu -> aurait pu, typo

Je présume que $C \circ C'$ est lui aussi muni de
\[
\mathrm{id}_{\Cat/S} \longrightarrow C \circ C',
\]

\msp{page 197}
et même a priori de deux façons possible
\[
\begin{tikzcd}[column sep=large, row sep=large]
& C(C'(X)) & \\
C(X) \arrow[ur, "C(i'_X)"] & & C'(X) \arrow[ul, "i_{C'(X)}"'] \\
& X, \arrow[ul, "i_X"] \arrow[ur, "i'_X"'] \arrow[uu, dotted] &
\end{tikzcd}
\]
soit par $C(i'_X) \circ i_X$, soit par $i_{C(X)} \circ i'_X$, heureusement elles sont égales, car pour toute flèche $X \xrightarrow{\ i'\ } X'$ dans $\Cat/S$, le carré % [sic] i_{C(X)} -> i_{C'(X)}, the diagram above has i_{C'(X)}

% ---- from Der-ChXIIc-30-03-03_p016-020.tex
\printedpage{130}
% head: A. Grothendieck : Dérivateurs, Ch. XII, Édition des Universités de Montpellier II et Paris VII

\[
\begin{tikzcd}[column sep=small, row sep=small]
 & C(X') & \\
C(X) \arrow[ur, "C(i')"] & & X' \arrow[ul, "i_{X'}"'] \\
 & X \arrow[ul, "i_X"] \arrow[ur, "i'"'] &
\end{tikzcd}
\]

est commutatif (par fonctorialité de $i_X : X \longrightarrow C(X)$), et on applique ceci au cas de $X' = C'(X)$, $i' = i'_X$.

Cette remarque nous encourage à faire un saut dans l'abstraction, et d'introduire la catégorie\msp{page 198}
\[
\mathcal{C} = \mathcal{C}_S = \begin{cases}
\text{catégorie des couples } (C,i), \text{ avec} \\
C : \Cat/S \longrightarrow \Cat/S \text{ se factorisant par } \Fib_W S, \\
i : \id_{\Cat/S} \longrightarrow C \text{ tel que } i_X \in W \ \forall\ X \text{ dans } \Cat/S,
\end{cases}
\tag*{$(*)$}
\]
une flèche
\[
(C,i) \longrightarrow (C',i')
\]
étant une flèche
\[
C \xrightarrow{u} C' \qquad \text{(dans } \Hom(\Cat/S, \Cat/S))
\]
rendant commutatif le carré \texttt{[lire : le diagramme]}
\[
\begin{tikzcd}[column sep=small]
 & \id_{\Cat/S} \arrow[dl, "i"'] \arrow[dr, "i'"] & \\
C \arrow[rr, "u"] & & C'.
\end{tikzcd}
\]

En d'autres termes, $\mathcal{C}$ est défini comme une sous-catégiorie pleine de % [sic] sous-catégiorie -> sous-catégorie, misspelled word
\[
\id_{\Cat/S}\backslash\Hom(\Cat/S, \Cat/S),
\]
à savoir celle dont les objets $(C,i)$ satisfont les deux conditions $(*)$ plus haut (impliquant l'une et l'autre le localiseur $W$).

\msp{page 199}Ces conditions impliquent ceci

\begin{enumerate}
\item[a)] $\forall\ X$ dans Cat/$S$, $(u_X : C(X) \longrightarrow C'(X)) \in W$, puisque
\[
\begin{tikzcd}[column sep=small]
 & X \arrow[dl, "i_X\in W"'] \arrow[dr, "i'_X\in W"] & \\
C(X) \arrow[rr, "u_X"] & & C'(X),
\end{tikzcd}
\]
\printedpage{131}
% head: A. Grothendieck : Dérivateurs, Ch. XII, Édition des Universités de Montpellier II et Paris VII
et

\item[b)] si $W$ satisfait W(4b), de sorte que $u_X \in W^{\mathrm{f}}_S$, alors ce définit un \textit{isomorphisme} entre $\varphi_C(X)$ et $\varphi_{C'}(X)$ % [sic] ce définit -> ceci définit, missing word
\[
\begin{tikzcd}[row sep=small]
s \arrow[r, mapsto] & \hot_W(C(X)_s) \arrow[d, "\wr\, u_{X,s}"] \\
s \arrow[r, mapsto] & \hot_W(C'(X)_s).
\end{tikzcd}
\]
\end{enumerate}

Mieux encore, dans ce cas, le foncteur composé
\[
\mathcal{C} \times \Cat/S \longrightarrow \Fib_W S \longrightarrow \HOT^{\mathrm{lc}}_W(S)
\]
donne naissance à
\[
\mathcal{C} \longrightarrow \Hom(\Cat/S, \HOT^{\mathrm{lc}}_W(S))
\]
qui transforme toute flèche de $\mathcal{C}$ en \texttt{[un]} \textit{isomorphisme,} donc se factorise par le groupoïde fondamental de $\mathcal{C}$,
\[
\Pi\mathcal{C} \longrightarrow \Hom(\Cat/S, \HOT^{\mathrm{lc}}_W(S)).
\]

\msp{page 200}Quant au foncteur envisagé tantôt
\[
\begin{aligned}
\mathcal{C} &\quad\longrightarrow\quad \underbrace{\Hom\,(\Cat/S, \Hom(S^\circ, \Hot_W))}_{\simeq\Hom((\Cat/S)\times S^\circ, \Hot_W)} \\
C &\quad\longmapsto\quad ((X,s) \longmapsto \hot_W(C(X)_s))\,,
\end{aligned}
\]
il se déduit du précédent en composant avec
\[
\begin{tikzcd}[column sep=small]
\HOT_W(S) \arrow[rr] \arrow[dr, dotted] & & \Hom(S^\circ, \Hot_W), \\
 & \Hom(\Pi S^\circ, \Hot_W) \arrow[ur, dotted] &
\end{tikzcd}
\]
plus exactement, avec le foncteur qui s'en déduit
\[
\Hom(\Cat/S, \HOT_W(S)) \longrightarrow \Hom(\Cat/S, \Hom(S^\circ, \Hot_W)).
\]

On voit poindre ici un principe plausible de canonicité d'isomorphismes, en ce qui concerne les foncteurs
\[
C_W : \Cat/S \longrightarrow \HOT^{\mathrm{lc}}_W(S)
\]
définis par les $C \in \mathrm{Ob}\,\mathcal{C}$, par la relation
\[
\Pi\mathcal{C} \approx \ \text{catégorie ponctuelle, i.e. } \mathcal{C} \text{ est 1-connexe.}
\]

De plus, les arguments utilisés jusqu'ici sont de nature à tel point formelle, que les aspects spécifiques du Modéliseur Cat
\printedpage{132}
% head: A. Grothendieck : Dérivateurs, Ch. XII, Édition des Universités de Montpellier II et Paris VII
\msp{page 201}semblent n'y jouer pratiquement aucun rôle, ou un rôle très effacé.

Le carré de la page 197 montre surtout que la catégorie $\mathcal{C}$ est filtrante - $(C,i)$, $(C,i')$ s'envoient l'un et l'autre dans $(C'',i'')$, avec $C'' = C \circ C'$, et $i'' = (C * i') \circ i = (i * C') \circ i'$. Il y a même une loi de composition dans $\mathcal{C}$, nullement commutative ni unitaire c'est sûr, mais sûrement associative, et (si cette composition est notée $*$) avec des morphismes fonctoriels % [sic] (C,i') -> (C',i'), the second couple of the following formula is (C',i')
\[
\xi \xrightarrow{\alpha_{\xi,\xi'}} \xi * \xi' \xleftarrow{\beta_{\xi,\xi'}} \xi'
\]
qui assurent le caractère filtrant croissant de $\mathcal{C}$. Ces flèches semblent remplacer en quelque sorte l'unité bilatère absente, et on s'attend à ce qu'elles aient les mêmes excellentes propriétés que si on avait une telle unité, concernant les composés doubles $(\xi * \xi') * \xi'' \simeq \xi * (\xi' * \xi'')$,

\msp{page 202}
\[
(\xi * \xi') * \xi'' \xrightarrow{\sim} \xi * (\xi' * \xi'').
\]

Il y a deux façons d'envoyer un produit binaire
\[
\xi * \xi', \qquad \xi' * \xi'', \qquad \xi * \xi''
\]
dans le produit ternaire, et bien elles doivent dans chacun des trois cas, être égales. Et de même il y a, ceci admis, encore deux façons d'envoyer $\xi$ dans le produit ternaire, en factorisant soit par $\xi * \xi'$, soit par $\xi * \xi''$, et de même pour les deux autres foncteurs $\xi'$, $\xi''$; et là encore, ces deux façons doivent être égales, dans chacun de ces trois cas - donc $3 + 3 = 6$ égalités en tout à écrire et à vérifier - on a vu pire!

La philosophie calculatoire devrait être celle-ci : Si on a un produit $n$-uple
\[
\xi_1 * \xi_2 * \cdots * \xi_n
\]
(où je me dispense de mettre des parenthèses \dots), pour tout produit $m$-uple partiel $(m \in [1, n-1])$\msp{page 203}
\[
\xi_{i_1} * \xi_{i_2} * \cdots * \xi_{i_m} \qquad\qquad (i_1 < i_2 < \cdots < i_m),
\]
en utilisant à tour de bras les isomorphismes d'associativité et les $\alpha_{\xi,\xi'}$, $\beta_{\xi,\xi'}$, il n'y a pourtant qu'une seule façon de l'envoyer dans le produit $n$-uple. J'ai bien l'impression que ça doit être vrai. (Est-ce que ça implique déjà que $\mathcal{C}$ est 1-connexe?)

En fait, il y a bien une unité bilatère, mais elle n'est pas dans $\mathcal{C}$, mais dans la catégorie plus grande
\[
\id_{\Cat/S}\backslash\Hom(\Cat/S, \Cat/S),
\]
qui elle, comme toute catégorie
\[
\id_{\mathcal{X}}\backslash\Hom(\mathcal{X}, \mathcal{X}),
\]
\printedpage{133}
% head: A. Grothendieck : Dérivateurs, Ch. XII, Édition des Universités de Montpellier II et Paris VII
est une catégorie monoïdale, i.e. avec produit associatif \textit{et} unitaire. Les flèches $\alpha_{\xi,\xi'}$, $\beta_{\xi,\xi'}$ sont les flèches qui résultent de la présence de cette unité (vu que $\mathcal{C}$ est une sous-catégorie pleine), et alors il doit devenir clair qu'on a les compatibilités dites.

\msp{page 204}Mais même une catégorie monoïdale, est elle 1-connexe? Par exemple $B_G$, pour $G$ un groupe disons, est bien une catégorie monoïdale et même ‘groupale’, une $G$-catégorie (tout ce qu'il y a de strict), et pourtant elle n'est \textit{pas} 1-connexe!

Bien sûr, on peut court-circuiter la question de la 1-connexité de $\mathcal{C}$, en notant qu'après tout on a déjà a priori un objet privilégié dans
\[
\Hom(\Cat/S, \HOT^{\mathrm{lc}}_W(S)),
\]
i.e. un foncteur canonique
\[
\Cat/S \longrightarrow \HOT^{\mathrm{lc}}_W(S),
\]
savoir le foncteur composé
\[
\Cat/S \underbrace{\longrightarrow}_{\substack{\text{localisation}\\ \text{par } W_S}} (\Cat/S)W_S^{-1} \longrightarrow \underbrace{\HOT^{\mathrm{lc}}_W(S)}_{\overset{\text{déf}}{=} (\Fib_W S)(W^{\mathrm{f}}_S)^{-1}} \quad ,
\tag*{$(*)$}
\]
où la deuxième flèche est définie comme le quasi-inverse de
\[
(\Fib_W S)(W^{\mathrm{f}}_S)^{-1} \longrightarrow (\Cat/S)W_S^{-1}
\tag*{$(**)$}
\]

\msp{page 205}déduit de l'inclusion
\[
(\Fib_W S, W^{\mathrm{f}}_S) \hookrightarrow (\Cat/S, W_S).
\]

Le fait que $(*)$ soit une équivalence, se prouve justement par le fait de l'existence d'un foncteur $C : \Cat/S \longrightarrow \Cat/S$ et d'un $i : \id_{\Cat/S} \longrightarrow C$, ayant les propriétés qu'on sait, i.e. du fait que $C \neq \emptyset$. De plus, la donnée d'un objet $C$ de $\mathcal{C}$ donne un moyen commode d'expliciter le quasi-inverse dudit foncteur. Mais ce quasi-inverse a une existence canonique et indépendante du choix d'un $C$, i.e. d'un mode de calcul. Le fait que le foncteur composé $\varphi$ dans $(*)$ puisse se ‘relever’ (à isomorphisme près) en un foncteur % [sic] C \neq \emptyset -> \mathcal{C} \neq \emptyset, the non-emptiness concerns the category \mathcal{C}
\[
C : \Cat/S \longrightarrow \Fib_W S
\]
entre $\varphi(X) = \hot_{W,S}(X)$ et $\hot_{W,S}(C(X))$ et que l'isomorphisme puisse s'expliciter, lui aussi, par une flèche $X$-fonctorielle
\[
i_X : X \longrightarrow C(X), \qquad i_X \in W,
\]
dans Cat/$S$ lui-même (et non seulement dans le localisé $(\Cat/S)W_S^{-1}$),
\printedpage{134}
% head: A. Grothendieck : Dérivateurs, Ch. XII, Édition des Universités de Montpellier II et Paris VII
\msp{page 206}est une vertu particulière du \texttt{[mot illisible]}, qui ne doit pas pour autant faire oublier le caractère intrinsèque des flèches $(*)$, indépendamment de tout mode de calcul.

\marginpar{$\|$}% double vertical bar printed in the left margin beside the first lines of this paragraph
La clef de la situation, et des questions de canonicité, me semble être dans le fait que $(**)$ est une équivalence de catégories - et il faut avouer que, tout familier et quasiment banal qu'il soit finalement devenu pour moi, c'est là un fait tout à fait remarquable et nullement évident a priori. (Mais sans doute peut-on l'avoir aussi par les fourbis généraux à la \textsc{Quillen-Thomason}, sans se battre avec un formalisme de chemins \dots) Les ‘isomorphismes canoniques’ qu'on peut établir entre une ‘théorie $C$ ’ et une ‘théorie $C'$ ’ (via deux constructions différentes, p. ex., de catégories de chemins, ou via l'approche de \textsc{Quillen-Thomason}, sans chemins) doivent toujours revenir à ceux

\msp{page 207}qu'on peu définir via la comparaison de la ‘théorie $C$ ’ d'une part, la ‘théorie $C'$ ’ de l'autre, avec ce que nous donne directement la flèche composée $(*)$ (p. 204). Finalement, il me semble que \textit{l'existence} d'une théorie des ‘fibres homotopiques’ (ou ‘$W$-homotopiques’) d'une flèche $X \longrightarrow S$ dans Cat, résulte de l'équivalence $(**)$ (p. 204), et cette théorie pourrait sans doute dans une certaine mesure se développer à partir de là, sans choix d'un foncteur $C$ ni d'une structure de catégorie de modèles au sens de \textsc{Quillen}. % [sic] qu'on peu définir -> qu'on peut définir, verb form

\section*{4\quad Axiomatisation des ‘fibres $W$-homotopiques’.}

La reflexion de hier me fait revenir sur l'axiomatique dans des ‘catégories de modèles’ % [sic] reflexion -> réflexion, missing accent
\[
(\mathcal{M}, W)
\]
(\footnote[119]{$W$ satisfait l'axiome de saturation faible : les isomorphismes \texttt{[sont]} dans $W$, et si deux parmi $f$, $g$, $gf$ \texttt{[sont]} dans $W$, le troisième aussi.}) plus générales que Cat, dans l'esprit de VI (cf. théorème-scholie 2, p. 108-113). Pour le moment, on va supposer seulement $\mathcal{M}$ stable par produits fibrés (on n'a pas besoin d'objet final). On introduit l'ensemble
\[
\mathrm{Fib}_W \subseteq \mathrm{Fl}(\mathcal{M})
\]

\msp{page 208}des \textit{$W$-fibrations} (sous-entendu ‘ultrafortes’!) par la condition

\begin{quote}
$f : X \longrightarrow Y$ est dans $\mathrm{Fib}_W$ si et seulement si le foncteur changement de base par $X$,
\[
\begin{aligned}
Y' &\quad\longmapsto\quad X' = X \times_Y Y' \\
\mathcal{M}/Y &\quad\longrightarrow\quad \mathcal{M}/X
\end{aligned}
\]
\end{quote}

% ---- from Der-ChXIIc-30-03-03_p021-025.tex
\printedpage{135}
% head: A. Grothendieck : Dérivateurs, Ch. XII, Édition des Universités de Montpellier II et Paris VII

\begin{quote}
transforme $W_X$ dans $W_Y$, où $W_X \subseteq \Fl(\cM/X)$ \texttt{[est]} défini par

$u \in W_X \iff \sigma(U) \in W$ \ ($\sigma : \cM/X \longrightarrow \cM$ le foncteur ‘source’, ou ‘oubli de $X$’). % [sic] \sigma(U) -> \sigma(u), the arrow is called u
\end{quote}

(\footnote[120]{\textbf{NB} Donner ici la proposition: $\mathrm{Fib}_W$ \textit{stable par composition et par changement de base, et contient les isomorphismes, stable par facteurs directs.}}). Posons, pour tout $S$ dans $\cM$,
\[
\Ho^!_W(S) = (\cM/S)W_S^{-1},
\]
qui joue le rôle d'une catégorie de $W$-types d'homotopie relatifs \textit{localement constants} sur $S$. On aimerait décrire une loi contravariante en $S$ pour $\Ho^!_W(S)$, i.e. pour
\[
f : S' \longrightarrow S
\]
flèche quelconque dans $\cM$, définir
\[
f^*_W : \Ho^!_W(S) \xrightarrow{?} \Ho^!_W(S'). \tag*{$(*)$}
\]
La difficulté est, bien sûr, que le foncteur changement de base
\[
f^* : \cM/S \longrightarrow \cM/S'
\]
n'est \textit{pas} en général compatible avec les \msp{page 209} localiseurs $W_S$, $W_{S'}$. La notion de $W$-fibration sert donc en premier lieu à dégager la condition la plus évidente sur une flèche $f : S' \longrightarrow S$ dans $\cM$, pour que le foncteur changement de base soit compatible avec les localiseurs $W_S$, $W_{S'}$ (c'est \textit{ça} la définition), donc induit $(*)$ sur les localisées.

Le \textit{‘problème des fibres homotopiques’} est justement de définir $(*)$ pour \textit{toute} flèche dans $\cM$, pas seulement pour les $W$-fibrations, de tellle façon que les $\Ho^!_W(S)$ forment, pour $S$ % [sic] tellle -> telle, letter doubled
variable, les fibres d'une catégorie fibrée remarquable sur $\cM$. Quand $\cM$ admet un objet final, de sorte que l'on peut regarder les flèches
\[
p : e \longrightarrow S \qquad (p \text{ comme ‘point’})
\]
dans $\cM$ comme des ‘points’ de $S$, la ‘fibre $W$-homotopique’ d'un objet quelconque $X/S$ de $\cM/S$ en ledit point sera par définition l'élément de
\[
\Ho_W = \Ho_W(\cM) \simeq \Ho^!_W(e)
\]
défini par
\[
p^*(\ho^!_{W/S}(X/S)),
\]
\printedpage{136}
% head: A. Grothendieck : Dérivateurs, Ch. XII, Édition des Universités de Montpellier II et Paris VII
\msp{page 210}
où
\[
\ho^!_{W/S} : \cM/S \longrightarrow \Ho^!_W(S) \overset{\text{déf}}{=} (\cM/S)(W_S)^{-1}
\]
est le foncteur de localisation canonique.

Pour tout $S$ dans $\cM$, soit
\[
\Fib_W S \subseteq \cM/S
\]
la sous-catégorie pleine formée des objets $X/S$ de $\cM/S$ tels que $X \longrightarrow S$ soit $W$-fibrant. On va définir une variante
\[
\Ho^{\lc}_W(S)
\]
de $\Ho^!_W(S)$, comme
\[
\Ho^{\lc}_W(S) \overset{\text{déf}}{=} (\Fib_W S)(\Wu_S)^{-1},
\]
où

\begin{quote}
$\Wu_S \subseteq \Fl(\Fib_W S)$ est le localiseur formé des flèches $u : X \longrightarrow Y$ qui sont ‘universellement sur $S$ dans $W$’, i.e. telles que pour toute flèche $S' \longrightarrow S$ de $\cM$ de but $S$, la flèche déduite par changement de base,
\[
u' : X' = X \times_S S' \longrightarrow Y' = Y \times_S S'
\]
est dans $W_{S'}$.
\end{quote}

On a ‘fait ce qu'il fallait’ pour que $\Ho^{\lc}_W(S)$ soit contravariant en $S$ pour des flèches \textit{quelconques} $f : S' \longrightarrow S$ de $\cM$, elles définissent
\[
f^{*\,\lc}_W : \Ho^{\lc}_W(S) \longrightarrow \Ho^{\lc}_W(S'), \tag*{$(**)$}
\]
\msp{page 211}
puisque les $W$-fibrations sont stables par changement de base, et que par suite $f^* : \cM/S \longrightarrow \cM/S'$ induit
\[
f^*_{\Fib_W} : \Fib_W S \longrightarrow \Fib_W S'
\]
qui, par définition de $\Wu_S$ et $\Wu_{S'}$, est compatible avec les localiseurs
\[
f^*_{\Fib_W} : (\Fib_W S, \Wu_S) \longrightarrow (\Fib_W S', \Wu_{S'}),
\]
et passe donc aux catégories des fractions en une flèche $(*)$.

On peut se demander pourquoi ne pas être plus généreux, et définir $\Ho^{\lc}_W(S)$ comme $(\cM/S)(\Wu_S)^{-1}$, en définissant $\Wu_S$ dans $\cM/S$ de la même façon que tantôt dans $\Fib_W S$ - ce qui ne coûte pas plus cher. Une première réponse est d'expérience le localiseur $\Wu_S$ dans $\cM/S$ tout entier s'avère pratiquement inutilisable, d'une part parce que l'appartenance à $\Wu_S$ (pour une flèche entre objets de $\cM/S$ qui ne sont pas très particuliers, p.ex. $W$-fibrés sur $S$) est invérifiable (sauf pour des isomorphismes et des choses comme ça \dots), et que d'autre part la catégorie de fractions est prohibitivement
\printedpage{137}
% head: A. Grothendieck : Dérivateurs, Ch. XII, Édition des Universités de Montpellier II et Paris VII
\msp{page 212}
grosse et ne peut servir (semble-t-il) à rien. Une autre réponse, c'est qu'on veut que la catégorie obtenue réponde bien à l'intuition de types d'homotopie relatifs \textit{localement constants} sur $S$ - et on veut donc que pour un changement de base quelconque $f : S' \longrightarrow S$ par un

\begin{quote}
$S'$ dans $\cM/S$, d'où un foncteur composé
\[
\begin{tikzcd}[column sep=small]
\cF_S(\Wu_S)^{-1} \arrow[r, "f^{*\,\lc}_W"] \arrow[rr, dotted, bend right=15, "\varphi_{S'}"'] & \cF_{S'}(\Wu_{S'})^{-1} \arrow[r] & \Ho_W
\end{tikzcd}
\]

(\footnote[121]{$\cF$ désigne ici une sous-catégorie fibrée de $\cM/\cM = \underline{\Fl}(\cM)$ sur $\cM$, substitut éventuel de celle formée de $\Fib_W S$ ou de $\cM/S$ tout entier.}) déduit des foncteurs de catégories avec localiseurs
\[
(\cF_S, \Wu_S) \xrightarrow{f^*} (\cF_{S'}, \Wu_{S'}) \xrightarrow{\text{oubli de } S'} (\cM, W)
\]
\end{quote}

(foncteur composé $\varphi_{S'}$ qui peut être vu comme le foncteur \textit{‘fibre homotopique en l'objet $S'$ sur $S$’}, cet objet étant visualisé comme un ‘point’ de $S$ ‘à valeurs dans $S'$ ’ suivant le yoga fonctoriel général \dots), le foncteur
\[
\begin{aligned}
S' &\longmapsto \varphi_{S'}(\xi) \qquad \text{pour } \xi \text{ dans } \cF_S(\Wu_S)^{-1}\\
\cM/S &\longrightarrow \Ho_W
\end{aligned}
\]
transforme flèches dans $W_S$ en isomorphismes : la ‘fibre $W$-homotopique’ en $S'/S$ d'un objet de $\cF_S(\Wu_S)^{-1}$ (vu comme $W$-type d'homotopie \textit{localement constant} sur $S$) doit être ‘la même’, si $S'$ remplacé par $S''$ avec $S'' \longrightarrow S'$ donné et dans $W_S$. C'est pour cela qu'il est raisonnable de prendre $\cF_S$ ‘assez petit’ (savoir $\Fib_W S$, qui est \msp{page 213} exactement ajusté pour cela), plutôt que de prendre $\cM/S$ tout entier.

On peut donc dire que la catégorie $\Fib_W S$, et le localiseur $\Wu_S$ ont été construits ad-hoc, le plus simplement possible, pour donner lieu à un localisé $\Ho^{\lc}_W(S)$ qui réponde à l'intuition de $W$-types d'homotopie relatifs localement constants des deux façons suivantes :

\begin{enumerate}[label={}]
\item[1°)] Il y a spontanément une loi contravariante en $S$, pour des flèches quelconques $S' \longrightarrow S$ dans Cat (on a défini $\Wu_S$ juste exprès en conséquence).
\item[2°)] que la notion de ‘fibre $W$-homotopique en $S'/S$’ d'un objet $\xi$ de $\Ho^{\lc}_W$ satisfasse à la condition précédente, i.e. une flèche $S'' \longrightarrow S'$ dans $\cM/S$ qui est dans $W_S$ définisse un \textit{isomorphisme} (dans $\Ho_W$) des fibres homotopiques
\[
\underbrace{\xi_{S''}}_{\varphi_{S''}(\xi)} \longrightarrow \underbrace{\xi_{S'}}_{\varphi_{S'}(\xi)}
\]
(et on a choisi $\cF_S = \Fib_W(S)$ juste exprès pour ça).
\end{enumerate}

\printedpage{138}
% head: A. Grothendieck : Dérivateurs, Ch. XII, Édition des Universités de Montpellier II et Paris VII

On a donc une excellente notion de fibres $W$-homotopiques pour des $X$ $W$-fibrés sur $S$ - cette fois excellente pour des changements de base \textit{quelconque} - et il n'y a aucun mérite à ça puisqu'on a tout défini exprès pour ça! % [sic] quelconque -> quelconques, agreement with "des changements"

\msp{page 214}
Pour le dire autrement, on a défini \textit{deux} notions de $W$-types d'homotopie relatifs sur $S$, l'une de façon à donner un tel type d'homotopie $\ho^!_{W,S}(X)$ pour \textit{tout} objet de $\cM/S$, mais il n'y a a priori une bonne notion de changement de base que pour les changements de base $S' \longrightarrow S$ qui sont $W$-fibrés, l'autre, qui correspond plus directement à l'intuition d'un $W$-type d'homotopie localement constant sur $S$, n'associe un tel type $\ho^{\lc}_{W,S}(X)$ qu'aux $X$ de $\cM/S$ qui sont $W$-fibrés sur $S$, en revanche cette notion admet une bonne notion d'image inverse et de ‘fibre $W$-homotopique’ pour des changements de base quelconques.

Une ‘bonne’ théorie des fibres $W$-homotopiques doit réunir en une seule les deux notions précédentes, de sorte que les avantages des deux s'ajoutent, et les désavantages s'éliminent! Elle s'obtient à l'aide d'un axiome draconien sur la situation (alors que jusqu'à maintenant on n'a rien utilisé sur $W \subseteq \Fl(\cM)$, si ce n'est tout au plus la saturation faible, en plus de l'existence des produits fibrés) :

\msp{page 215}
\circled{a} (\textit{Axiome de la fibre homotopique.}) Pour tout objet $S$ dans $\cM$, le foncteur canonique
\[
\Ho^{\lc}_W(S) \longrightarrow \Ho^!_W(S)
\]
déduit de l'inclusion
\[
(\Fib_W(S), \Wu_S) \longrightarrow (\cM/S, W_S)
\]
est une \textit{équivalence de catégories.}

Cet axiome implique que dans $\Fib_W S$, $\Wu_S$ et le localiseur induit par $W_S$ (que je désigne encore par $W_S$ par abus de notation) ont même saturé fort. En fait, il est naturel de compléter l'axiome précédent par le suivant.

\circled{b} (Axiome adjoint de l'axiome de la fibre homotopique.) Sur $\Fib_W S$, les localiseurs $\Wu_S$ et $W_S$ coïncident, i.e. on a ceci :

\begin{quote}
Soient $X$, $Y$ dans $\Fib_W S$, $X \xrightarrow{u} Y$ un $S$-morphisme tel que $u \in W_S$, alors $u \in \Wu_S$ (i.e. le morphisme $u' : X' \longrightarrow Y'$ déduit par tout changement de base $S' \longrightarrow S$ est encore dans $W$).
\end{quote}

\msp{page 216}
Je présume que a) sera rarement satisfait sans que b) le soit (et inversement d'ailleurs), et que la démonstration de a) et b) se fera pratiquement en une seule haleine.

Si $\cM = \mathrm{Cat}$, $W$ un localiseur fondamental, alors la notion adaptée ici de $W$-fibration (lire : ‘ultraforte’, dans la terminologie spécifique développée dans Cat) est à tel point forte, qu'elle implique déjà pratiquement b) (moyennant W(4a$'$), il est vrai, cf. p. 53), et
\printedpage{139}
% head: A. Grothendieck : Dérivateurs, Ch. XII, Édition des Universités de Montpellier II et Paris VII
a) n'est guère vérifiable que si on a W(4a) (\footnote[122]{Non, il me semble que W(4b) doit suffire \dots}) (loc. cit.). Donc ici c'est ‘l'axiome des fibrations’ sur $W$, sous sa forme forte W(4a), qui est la clef de voûte d'une bonne théorie des fibres homotopiques.

Une telle théorie (revenant au cas général) revient à dire, pour l'essentiel, qu'il y a ‘assez’ de $W$-fibrations (‘ultrafortes’) pour pouvoir décrire \textit{tout} objet d'un $\cM/S$, à $W_S$-équivalence près, comme un objet $W$-fibrant sur $S$. Dans les cas à ma connaissance, on a de plus ceci : on \msp{page 217} a pour tout $S$ dans $\cM$ une sous-catégorie $\cF_{0\,S}$ de $\Fib_W S$ (je pense par exemple à $\underline{\mathrm{Parf}}_W S$, voire $\underline{\mathrm{Parf}}_{W_\omega} S$, dans le cas où $\cM = \mathrm{Cat}$), telle que pour toute sous-catégorie pleine $\cF$ de $\cM/S$ telle que
\[
\cF_{0\,S} \subseteq \cF \subseteq \cM/S
\]
(a fortiori, pour une $\cF$ pleine telle que $\Fib_W S \subseteq \cF \subseteq \cM/S$), les inclusions induisent des équivalences de catégories
\[
\cF_{0\,S}(W_S)^{-1} \xrightarrow{\sim} \cF(W_S)^{-1} \xrightarrow{\sim} \underbrace{(\cM/S)(W_S)^{-1}}_{\overset{\text{déf}}{=} \mathrm{Hot}^!_W(S)}. % [sic] Hot -> Ho, elsewhere the notation is Ho^!_W(S)
\]
Le choix de $\cF_{0\,S}$ ne semble pas pouvoir se fixer de façon ‘la meilleure possible’, une fois pour toutes dans le contexte axiomatique adopté ici. C'est un élément de structure supplémentaire dans $\cM$ (en plus du localiseur $W$), au même titre que la donnée des classes de $Q$-fibrations et $Q$-cofibrations, à la \textsc{Quillen}. J'ignore si les axiomes a), b) ci-dessus (axiome de la fibre homotopique) impliquent déjà \msp{page 218} la variante de \circled{a} :

\circled{c} Pour toute sous-catégorie pleine $\cF$ de $\cM/S$ telle que
\[
\Fib_W S \subseteq \cF \subseteq \cM/S
\]
les inclusions induisent des équivalences
\[
\underbrace{(\Fib_W S)(W_S)^{-1}}_{\begin{subarray}{l} = \ \Ho^{\lc}_W(S)\\ \text{moyennant a)}\end{subarray}} \xrightarrow{\sim} \cF(W_S)^{-1} \xrightarrow{\sim} \underbrace{(\cM/S)(W_S)^{-1}}_{=\Ho^!_W(S)}.
\]

Mais je présume que pratiquement, quand on pourra vérifier \circled{a}, on pourra vérifier \circled{c} également (\footnote[123]{Cf. page 240 ff. Il faudrait ici développer un peu plus le formalisme de $\Ho^{\lc}_W(S)$ et les fibres homotopiques, sous l'hypothèse \circled{a}\,\circled{b} \ plus haut.}).

% ---- from Der-ChXIIc-30-03-03_p026-030.tex
\printedpage{140}
% head: A. Grothendieck : Dérivateurs, Ch. XII, Édition des Universités de Montpellier II et Paris VII

\section*{5 Axiomatique des $W$-types d'homotopie relatifs non localement constants $\mathrm{Ho}^{\mathcal{L}}_W(S)$.}

On n'a rien dit jusqu'à ce point d'éventuels $W$-types d'homotopie relatifs (sur un objet $S$ de $\mathcal{M}$) qui ne seraient \textit{pas} ``localement constants''. Il en était question pourtant dans loc. cit. (VI, théorème-scholie 2, partie II du théorème pages 109-113.) Je profite de l'occasion pour revenir là dessus. Il faut introduire, dans chaque $\mathcal{M}/S$, un deuxième localiseur
\[
W^0_S \subseteq \mathrm{Fl}(\mathcal{M}/S), \qquad W^0_S \subseteq W_S
\]
(plus fin que $W_S$), ce qui permettra de définir
\msp{page 219}
\[
\mathrm{Ho}_W(S) \deff (\mathcal{M}/S)(W^0_S)^{-1}.
\]
Qu'il n'y ait pas de choix \textit{canonique} qui s'impose ici, déduit de la seule structure $(\mathcal{M},W)$, est claire sur l'exemple de $(\mathrm{Cat},W)$, auquel cas on trouve \textit{deux} choix symétriques, transformés l'un en l'autre par l'involution
\[
X \longmapsto X^\circ : (\mathrm{Cat},W) \xrightarrow{\sim} (\mathrm{Cat},W),
\]
à savoir $W^{\mathrm{g}}_S$ et $W^{\mathrm{d}}_S$. Le choix de l'un ou de l'autre, disons (pour fixer les idées) de $W^{\mathrm{g}}_S$, est associé à une notion d'\textit{isomorphisme local} (ou de `coisomorphisme local' dans le cas de $W^{\mathrm{d}}_S$) que je voudrais axiomatiser. Il semble bien que pour les principales catégories de modèles pour Hot (et pour ses variantes), on dispose de façon assez naturelle d'une telle notion d'isomorphisme local (sans l'embarras, le plus souvent, d'un automorphisme involutif fournissant une notion symétrique \dots).

On supposera donc, comme dans loc. cit. p. 94-96, donné une classe
\[
\mathcal{L} \subseteq \mathrm{Fl}(\mathcal{M})
\]
de morphismes, jounant le rôle de `flèches de % [sic] jounant -> jouant, faute de frappe
\msp{page 220}
localisation' ou `d'isomorphismes locaux', satisfaisant les conditions suivantes.

\begin{enumerate}
\item[$(*)$] $\mathcal{L}$ contient les isomorphismes, est stable par composition, et par changement de base.
\end{enumerate}

On pose, pour $S$ dans $\mathcal{M}$,
\[
\mathcal{L}_S \subseteq \mathcal{M}/S \qquad \text{sous-catégorie pleine des } S'/S \text{ telle que } S' \longrightarrow S \text{ \texttt{[soit]} dans } \mathcal{L}.
\]
On définit alors pour tout $S$ dans $\mathcal{M}$
\[
\begin{aligned}
&W^{\mathcal{L}}_S \subseteq \mathrm{Fl}(\mathcal{M}/S)\\
&(X \xrightarrow{u} Y) \in W^{\mathcal{L}}_S \iff \forall\, S' \in \mathcal{L}_S,\ (u_{S'} : X \times_S S' \longrightarrow Y \times_S S') \text{ est dans } W.
\end{aligned}
\]
% foot: Version 30 mars 2003

\printedpage{141}
% head: A. Grothendieck : Dérivateurs, Ch. XII, Édition des Universités de Montpellier II et Paris VII

C'est donc une variante du localiseur $W^{\mathrm{u}}_S$, où maintenant on ne prend comme changements de base que les flèches `$\mathcal{L}$-localisantes', i.e. $\in \mathcal{L}$. On a donc
\[
W^{\mathrm{u}}_S \subseteq W^{\mathcal{L}}_S \subseteq W_S
\]
en tant que parties de $\mathrm{Fl}(\mathcal{M}/S)$, et de même dans toute sous-catégorie pleine $\mathcal{F}$ de $\mathcal{M}/S$. Dans le cas de $\underline{\mathrm{Fib}}_W S$, moyennant l'axiome \circled{b} ci-dessus ($W^{\mathrm{u}}_S = W_S$ dans $\underline{\mathrm{Fib}}_W S$), les trois localiseurs coïncident.

On pose
\[
\mathrm{Ho}^{\mathcal{L}}_W(S) \text{ (ou } \mathrm{Ho}_W(S)\text{)} \ \deff (\mathcal{M}/S)(W^{\mathcal{L}}_S)^{-1},
\]
et on a donné dans loc. cit. (partie II du théorème-scholie) un résultat intéressant sur cet invariant, pour $S$ fixé. Mon propos
\msp{page 221}
présent est de voir ce qui se passe pour $S$ variable. Visiblement, une notion de $W$-type d'homotopie relatif sur un objet quelconque $S$ de $\mathcal{M}$ n'est satisfaisante, que si elle permet des images inverses par des flèches \textit{quelconques}
\[
f : S' \longrightarrow S
\]
dans $\mathcal{M}$, qui doivent définir absolument
\[
f^{*\mathcal{L}}_W \text{ ou } f^*_W : \mathrm{Ho}^{\mathcal{L}}_W(S) \longrightarrow \mathrm{Ho}^{\mathcal{L}}_W(S').
\]
On est donc confronté à la même perplexité que tantôt - et ce n'est \textit{pas,} cette fois, l'introduction de $\underline{\mathrm{Fib}}_W(S)$ qui va sauver la situation, car
\[
(\underline{\mathrm{Fib}}_W S)(\underbrace{W^{\mathcal{L}}_S}_{\substack{= W_S \\ = W^{\mathrm{u}}_S}})^{-1} \longrightarrow (\mathcal{M}/S)(W^{\mathcal{L}}_S)^{-1}
\]
n'a aucune envie, cette fois, d'être une équivalence de catégories, mais bien plutôt une inclusion (pleinement fidèle) des $W$-types d'homotopie relatifs constants sur $S$, dans les $W$-types d'homotopie quelconques. Il faut visiblement un axiome qui remplace l'axiome `de la fibre homotopique'.

\msp{page 222}

Dans le cas de Cat, la classe $\mathcal{L}$ des isomorphismes locaux (disons) permet de définir, comme ci-dessus, le localiseur $W^{\mathcal{L}}_S = W^{\mathrm{g}}_S$ dans les $\mathrm{Cat}/S$, comme formé des $u : X \longrightarrow Y$ sur $S$ qui sont dans $W$ et le restent par tout changement de base $S' \longrightarrow S$ qui est dans $\mathcal{L}$. Mais il y a lieu de \textit{saturer} $\mathcal{L}$, en regardant la classe $\mathcal{L}^\wedge$ des flèches $f : S \longrightarrow S'$ de Cat telles que le changement de base
\[
f^* : \mathrm{Cat}/S \longrightarrow \mathrm{Cat}/S'
\]
envoie $W^{\mathrm{g}}_S$ dans $W^{\mathrm{g}}_{S'}$. Ce $\mathcal{L}^\wedge$ est tautologiquement stable par composition, et contient les isomorphismes (mieux, contient $\mathcal{L}$ = ensemble des isomorphismes locaux). Il n'est pas
% foot: Version 30 mars 2003
\printedpage{142}
% head: A. Grothendieck : Dérivateurs, Ch. XII, Édition des Universités de Montpellier II et Paris VII
clair, par la définition donnée ici, qu'il soit stable par changement de base - on va renforcer la définition en conséquence, en exigeant, pour que $f : S' \longrightarrow S$ soit dans $\mathcal{L}^\wedge$, que pour tout $T$ sur $S$, considérons $f_T : T' = S' \times_S T \longrightarrow T$, $f^*_T : \mathrm{Cat}/T \longrightarrow \mathrm{Cat}/T'$ soit compatible avec les localiseurs $W^{\mathrm{g}}_T$, $W^{\mathrm{g}}_{T'}$. Alors $\mathcal{L}^\wedge$ est (tout comme $\mathcal{L}$) stable par composition et par changement de base, et de plus il contient $\mathcal{L}$. De plus, pour tout $S$, on a $W^{\mathcal{L}}_S = W^{\mathcal{L}^\wedge}_S$. Ce petit sorite pourrait se faire bien sûr pour tout modéliseur $(\mathcal{M},W)$. Mais la chose intéressante ici, c'est que $\mathcal{L}^\wedge$ contient l'ensemble $\mathrm{Liss}_W$ des flèches $W$-lisses. (J'ignore s'il est égal.) On aurait donc pu
\msp{page 223}
aussi bien partir de l'ensemble de flèches $\mathrm{Liss}_W$, pour définir les localiseurs $W^{\mathrm{g}}_S$ comme
\[
W^{\mathrm{g}}_S = W^{\mathrm{Liss}_W}_S \qquad \text{au lieu de } W^{\mathrm{g}}_S = W^{\mathrm{Isoloc}}_S,
\]
ce qui est équivalent.

D'autre part, en travaillant avec les localisés $\mathrm{HOT}^\circ_W(S) = (\mathrm{Cat}/S)(W^{\mathrm{g}}_S)^{-1}$, le rôle de $\underline{\mathrm{Fib}}_W S$ pour la théorie des $\mathrm{HOT}^{\mathrm{lc}}_W(S)$ est joué ici par $\underline{\mathrm{Prop}}_W(S)$. On voit donc que dans la théorie des $W$-types d'homotopie relatifs non localement constants, mode $\mathrm{HOT}^\circ_W$, les ensembles de flèches $\mathrm{Liss}_W$ et $\mathrm{Prop}_W$ jouent un rôle en quelque sorte complémentaire : les flèches $f : S' \longrightarrow S$ de $\mathrm{Liss}_W$ sont (pratiquement) celles qui permettent directement de définir
\[
f^{*\circ}_W : \mathrm{HOT}^\circ_W(S) \longrightarrow \mathrm{HOT}^\circ_W(S')
\]
comme induit par
\[
f^* : (\mathrm{Cat}/S, W^{\mathrm{g}}_S) \longrightarrow (\mathrm{Cat}/S', W^{\mathrm{g}}_{S'}),
\]
tandis que les flèches de $\mathrm{Liss}_W$ permettent de définir pour tout $S$ dans Cat la sous-catégorie pleine
\[
\underline{\mathrm{Prop}}_W S \subseteq \mathrm{Cat}/S,
\]
donnant lieu à une équivalence de catégories
\[
\left\{
\begin{aligned}
&(\underline{\mathrm{Prop}}_W S)(W^{\mathrm{u}}_S)^{-1} \xrightarrow{\sim} \underbrace{(\mathrm{Cat}/S)(W^{\mathrm{g}}_S)^{-1}}_{= \mathrm{HOT}^0_W(S)} % ?? uncertain: exposant de HOT sous l'accolade -- 0 (zéro) ou \circ
\\
&\text{avec}\\
&W^{\mathrm{u}}_S = W^{\mathrm{g}}_S \text{ sur } \underline{\mathrm{Prop}}_W S
\end{aligned}
\right.
\]
(ce qui visiblement est le substitut annoncé de ``l'axiome de la fibre $W$-homotopique''). Cette
\msp{page 224}
équivalence, qui permet de décrire $\mathrm{Hot}^\circ_W(S)$ en termes de $\underline{\mathrm{Prop}}_W S \subseteq \mathrm{Cat}/S$ et d'un localiseur $W^{\mathrm{u}}_S$ \textit{stable par} \textbf{tout} \textit{changement de base,} permet d'autre part de définir la loi contravariante attendue pour $\mathrm{HOT}^\circ_W(S)$.
% foot: Version 30 mars 2003

\printedpage{143}
% head: A. Grothendieck : Dérivateurs, Ch. XII, Édition des Universités de Montpellier II et Paris VII

D'autre part, quand on échange le rôle de $\mathrm{Liss}_W$ et $\mathrm{Prop}_W$, on trouve la théorie complémentaire du $\mathrm{HOT}_W(S)$,
\[
\left\{
\begin{aligned}
&\mathrm{HOT}_W(S) \deff (\mathrm{Cat}/S)(W^{\mathrm{d}}_S)^{-1} \xleftarrow{\sim} (\underline{\mathrm{Liss}}_W S)(W^{\mathrm{u}}_S)^{-1},\\
&\text{sur } \underline{\mathrm{Liss}}_W S,\ W^{\mathrm{d}}_S = W^{\mathrm{u}}_S,
\end{aligned}
\right.
\]
la description de $\mathrm{HOT}_W(S)$ en termes de $\underline{\mathrm{Liss}}_W S$ et de $W^{\mathrm{u}}_S$ permet d'en définir le caractère contravariant en $S$, pour \textit{toute} flèche $f : S' \longrightarrow S$ dans Cat. Tandis que si on veut pouvoir passer aux catégories de fractions sur $\mathrm{Cat}/S$ et $\mathrm{Cat}/S'$ tout entier (via $W^{\mathrm{d}}_S$, $W^{\mathrm{d}}_{S'}$), il faut se limiter aux changements de base $W$-propres.

Du coup, me vient la perplexité si cette complémentaire, voire cette dualité, entre deux ensembles de flèches, les `lisses' et les `propres', ne serait pas un fait général qui se retrouvait tout au moins dans les catégories de modèles proche de la description de HOT ou d'un $\mathrm{Hot}_W$.

\msp{page 225}

Il se pourrait bien que dans $\boldsymbol{\Delta}^\wedge$ on ait une dualité parfaite entre ces deux espèces de notions (il faudra que j'examine la chose), peut-être même dans toute catégorie $\mathcal{A}^\wedge$ (mais dans $\boldsymbol{\Delta}^\wedge$ il y a un automorphisme involutif $K \longmapsto K^\circ$, correspondant au $X \longmapsto X^\circ$ dans Cat \dots). Et après tout, même dans le contexte des espaces topologiques, les notions de propreté d'une part, d'asphéricité locale d'autre part, pour une flèche $X \longrightarrow Y$, sont bien définies, et pourraient bien s'ajuster comme le font les notions correspondantes dans Cat, et (qui sait) donner lieu à deux versions complémentaires, ou `duales' l'une de l'autre, d'une notion de type d'homotopie relatif.

N'en ayant pas le cœur net, je suis un peu indécis comment présenter une description axiomatique de la théorie des $\mathrm{Ho}^{\mathcal{L}}_W(S)$, pour une catégorie de modèles générale $(\mathcal{M},W)$. Dois-je me donner d'emblée
\[
\mathrm{Liss}, \mathrm{Prop} \subseteq \mathrm{Fl}(\mathcal{M}),
\]
stables par composition, par changement de base, contenant les isomorphismes, et tels que (peut-être)
\[
\mathrm{Liss} \cap \mathrm{Prop} \subseteq \mathrm{Fib}_W,
\]
\msp{page 226}
mais surtout ils permettent de définir, pour $S$ de $\mathcal{M}$
\[
\left\{
\begin{array}{lcccl}
W^{\mathrm{g}}_S & \deff & W^{\mathrm{Liss}}_S & \subseteq & \mathcal{F}(\mathcal{M}/S)\\
W^{\mathrm{d}}_S & \deff & W^{\mathrm{Prop}}_S & \subseteq & \mathcal{F}(\mathcal{M}/S)
\end{array}
\right.
\]
avec les propriétés suivantes :

\begin{enumerate}
\item[{\circled{a}}] Si $f : S' \longrightarrow S$ est dans Prop (resp. Liss), alors $f^* : \mathcal{M}/S \longrightarrow \mathcal{M}/S'$ transforme $W^{\mathrm{d}}_S$ dans $W^{\mathrm{d}}_{S'}$ (resp. $W^{\mathrm{g}}_S$ dans $W^{\mathrm{g}}_{S'}$),
\end{enumerate}
% foot: Version 30 mars 2003
\printedpage{144}
% head: A. Grothendieck : Dérivateurs, Ch. XII, Édition des Universités de Montpellier II et Paris VII
\noindent ce qui est une pure tautologie, et ne concerne pas encore les relations de Prop avec Liss, mais par contre :

\begin{enumerate}
\item[{\circled{b}}]
\[
\left\{
\begin{aligned}
&(\underline{\mathrm{Liss}}S)(W^{\mathrm{u}}_S)^{-1} \xrightarrow{\sim} (\mathcal{M}/S)(W^{\mathrm{d}}_S)^{-1} \text{ \texttt{[est]} une équivalence}\\
&\text{et } W^{\mathrm{u}}_S = W^{\mathrm{d}}_S \text{ sur } \underline{\mathrm{Liss}}S,
\end{aligned}
\right.
\]
et dualement
\[
\left\{
\begin{aligned}
&(\underline{\mathrm{Prop}}S)(W^{\mathrm{u}}_S)^{-1} \xrightarrow{\sim} (\mathcal{M}/S)(W^{\mathrm{g}}_S)^{-1} \text{ \texttt{[est]} une équivalence}\\
&\text{et } W^{\mathrm{u}}_S = W^{\mathrm{g}}_S \text{ sur } \underline{\mathrm{Prop}}S.
\end{aligned}
\right.
\]
\end{enumerate}

De plus, il y a lieu d'énoncer \circled{c} deux propriétés d'adjonction (duales l'une de l'autre), pour lesquelles je renvoie à VI, théorème-scholie 2, p. 108-113. Il se pose ici la question subsidiaire de simplifier les conditions de validité de la formule d'adjonction, un peu hétéroclites dans loc. cit.

Il se pose d'autre part la question de faire la théorie des lois covariantes $f_!$ et $f_*$ pour $\mathrm{Ho}^{\mathrm{lc}}_W(S)$ et pour $\mathrm{Ho}^{\mathrm{Prop}}_W(S)$, $\mathrm{Ho}^{\mathrm{Liss}}_W(S)$.

\msp{page 227}

J'ai bon espoir que le foncteur $f_!$ peut être décrit, et son adjonction avec $f^*$ établie, à aussi bon compte dans ce contexte général que dans Cat. Mais je n'ai pas envie de m'y arrêter à présent.

\section*{6 Cofibre $W$-homotopique dans $\mathrm{Cat}$.}

Par contre, j'ai envie de regarder ici sous quelles conditions sur le localiseur fondamental $W$, il existe dans $(\mathrm{Cat},W)$ une théorie duale de celle de la fibre $W$-homotopique.

La question-clef est, bien sûr, celle d'avoir une bonne prise sur les $W$-cofibrations. Or je ne connais jusqu'à présent \textit{que} les immersions ouvertes de Dwyer, les immersions fermées de Dwyer, et ce qui s'en déduit par composition etc. Le travail de \textsc{Thomason} peut motiver l'espoir que la connaissance de ces immersions-là, qui semblent ici jouer un rôle similaire (dual) à celui des morphismes $W$-parfaits, suffira pour faire une bonne théorie des `cofibres $W$-homotopiques'.

% trois traits verticaux imprimés dans la marge gauche de ce paragraphe
\marginpar{\textbar\textbar\textbar}Je suppose dans toute cette section que $W$ satisfait W(7 bis) (axiome des carrés cocartésiens), pour avoir des bonnes propriétés des immersions ouvertes de Dwyer (\footnote[124]{\begin{enumerate}
\item[a)] Ce sont des $W$-cofibrations (résulte de b) et c)).
\item[b)] Si
\end{enumerate}}), et plus généralement, pour les $W$-cofibrations. % ?? uncertain: coupure en fin de ligne "co-cartésiens" -- cocartésiens ou co-cartésiens
% foot: Version 30 mars 2003

% ---- from Der-ChXIIc-30-03-03_p031-035.tex
\printedpage{145}
% head: A. Grothendieck : Dérivateurs, Ch. XII, Édition des Universités de Montpellier II et Paris VII
% foot: Version 30 mars 2003

\msp{page 228}
Il se pose donc la question de la factorisation duale d'une flèche $f : X \longrightarrow Y$ de Cat en
\[
\begin{tikzcd}
X \arrow[r, hook, "W\text{-cofibration}", "i_f"'] & \tilde{Y} \arrow[r, "\in W", "p_f"'] & Y.
\end{tikzcd}
\]

Les résultats de \textsc{Thomason} assurent que ça existe, avec $i$ une immersion ouverte de Dwyer, et $p$ une $W$-fibration - du moins dans le cas $W = W_\infty$. Je vais essayer de le voir directement, en calquant dualement la construction de \textsc{Cartan-Serre}.

Il s'agit donc d'abord de factoriser
\[
X \amalg X \xrightarrow{\ \mathrm{codiag}_X\ } X,
\]
et pour ceci je choisis un intervalle fixe
\[
I \ = \ \Phi_n \ = \ (0 \longrightarrow 1 \longleftarrow 2 \longrightarrow \cdots \longrightarrow 2n-1 \longleftarrow 2n)
\]
avec
\[
n \geq 2
\]
de sorte que l'inclusion
\[
\varepsilon = \text{catégorie discrète } \{0,1\} \xrightarrow{\partial} I, \qquad \partial(0) = 0, \quad \partial(1) = 2n,
\]
soit une immersion ouverte de Dwyer. Cette propriété est stable par changement de base Cat-fibrant, en particulier en prenant un produit
\[
\underbrace{X \times \varepsilon}_{\simeq X \coprod X} \xrightarrow{i_X} X \times I,
\]
% continuation of a footnote begun on the preceding page of the edition (no footnote mark on this page)
\begingroup\renewcommand{\thefootnote}{}\footnotetext{%
\[
\begin{tikzcd}[ampersand replacement=\&]
U \arrow[r, hook, "i"] \arrow[d, "p"'] \& X \arrow[d, "q"] \\
U' \arrow[r, hook, "i'"] \& X'
\end{tikzcd}
\]
avec $i$ immersion ouverte de Dwyer, alors
\[
\begin{aligned}
i \in W \quad &\Longleftrightarrow \quad i' \in W \\
p \in W \quad &\Longleftrightarrow \quad q \in W,
\end{aligned}
\]
j'ai donc besoin de W(7 bis) \textit{et} même de W(7) pour la deuxième implication.
\begin{enumerate}
\item[c)] Les immersions ouvertes de Dwyer \texttt{[sont]} stables par cochangement de base et par composition, contiennent les isomorphismes.
\end{enumerate}}\endgroup
\printedpage{146}
% head: A. Grothendieck : Dérivateurs, Ch. XII, Édition des Universités de Montpellier II et Paris VII
% foot: Version 30 mars 2003
\msp{page 229}
d'autre part
\[
X \times I \xrightarrow{\mathrm{pr}_1} X
\]
est $W_\omega$-parfait, d'où la factoriastion cherchée % [sic] factoriastion -> factorisation, misprint
\[
X \amalg X \xrightarrow{i_X} X \times I \ \xrightarrow{\ p_X = \mathrm{pr}_1^X\ } X
\]
- il n'y a pas plus simple! Comment l'utiliser pour construire la factorisation d'une flèche $f : X \longrightarrow Y$ quelconque, en dualisant les constructions relatives aux espaces de chemins? On définit $Z(f)$ par le carré cocartésien dans Cat
\[
\begin{tikzcd}
X \arrow[r, dotted, "\mathrm{inj}_1^X"'] & X \times I \arrow[r, "f'"] & Z(f) \arrow[d, dotted, bend left, "p_f"] \\
& X \arrow[u, hook, "\mathrm{inj}_0^X"'] \arrow[r, "f"] & Y, \arrow[u, hook, "W"']
\end{tikzcd}
\]
on a
\[
Z(f) \xrightarrow{\alpha} Y \qquad \text{induit par } \mathrm{pr}_1 : X \times I \longrightarrow X,
\]
et la factorisation cherchée de $f$ est
\[
X \ \xrightarrow{\ i_f = f' \circ \mathrm{inj}_1^X\ } Z(f) \xrightarrow{p_f} Y.
\]

Je dis que
\[
\left\{
\begin{aligned}
& i_f \text{ est immersion ouverte de Dwyer, donc une } W\text{-cofibration.} \\
& p_f \in W.
\end{aligned}
\right.
\]

\msp{page 230}
On voit sur le diagramme plus complet
\[
\begin{tikzcd}
& X \times I \arrow[r, "f'"] \arrow[dd, dotted, bend left=20, "\mathrm{pr}_1"] \arrow[dr, phantom, "\substack{\text{cocart.}\\ \text{et cart.}}"] & Z(f) \arrow[dd, dotted, bend left=60, "p_f"] \\
X \arrow[ur, dotted, "\mathrm{inj}_1^X"] \arrow[r, hook, "i_1^X"] & X \amalg X \arrow[u, hook, "i_X"'] \arrow[r, "f \coprod \mathrm{id}_X"] \arrow[dr, phantom, "\substack{\text{cocart.}\\ \text{et cart.}}"] & Y \amalg X \arrow[u, hook, "j"'] \\
& X \arrow[u, hook, "i_0^X"'] \arrow[r, "f"] & Y, \arrow[u, hook]
\end{tikzcd}
\]
\printedpage{147}
% head: A. Grothendieck : Dérivateurs, Ch. XII, Édition des Universités de Montpellier II et Paris VII
% foot: Version 30 mars 2003
\[
i_f = f' \circ \mathrm{inj}_1^X = f' i_X i_1^X = j \ \ \underbrace{(f \amalg \mathrm{id}_X) i_1^X}_{\substack{\text{inclusion canonique}\\ \text{de } X \text{ dans}\\ Y \coprod X}} \ \ ,
\]
c'est \texttt{[une]} immersion ouverte de Dwyer, comme composé de $Y \longrightarrow Y \amalg X$ et $j : Y \amalg X \longrightarrow Z(f)$ qui le sont ($j$ obtenu par cochangement de base à partir de $i_X$, qui est immersion ouverte de Dwyer).

D'autre part, $X \xrightarrow{\mathrm{inj}_0^X} X \times I$ est \texttt{[une]} immersion ouverte de Dwyer et dans $W$, donc coüniversellement dans $W$, donc $Y \longrightarrow Z(f)$ est elle aussi dans $W$, donc aussi la rétraction $p_f$.

D'autre part, la proposition duale de
\[
\underline{\mathrm{Fib}}_W S \cap W_S = W^{\mathrm{u}}_{S\,\mathrm{univ}}
\]
est valable pour les $W$-cofibrations :

\msp{page 231}
\textbf{Proposition 16} (\footnote[125]{N'a rien à voir avec Cat - une chose générale quand il y a assez de $W$-cofibrations.}). \textit{Sur la sous-catégorie pleine $\underline{\mathrm{Cofib}}_W(S)$ de $S\backslash\mathrm{Cat}$, le localiseur ${}_SW$ est égal au $({}_SW)^{\mathrm{univ}}$ ($W$-équivalences couniverselles), i.e. on a ceci : Soient}
\[
S \longrightarrow X, \qquad S \longrightarrow Y
\]
\textit{des $W$-cofibrations, et soit $f : X \longrightarrow Y$ un morphisme de $S\backslash\mathrm{Cat}$ tel que $f \in W$. Alors pour tout cochangement de base $S \longrightarrow S'$, la flèche correspondante}
\[
f' : X' = X \amalg_S S' \longrightarrow Y' = Y \amalg_S S'
\]
\textit{est dans $W$.}

\textsc{Démonstration}. On factorise $S \longrightarrow S'$ en
\[
S \xrightarrow{i} S'_0 \xrightarrow{j} S',
\]
avec $i$ immersion ouverte de Dwyer (donc $W$-cofibration), et $j \in W$. On a donc le diagramme à carrés cocartésiens
\[
\begin{tikzcd}
S \arrow[r] \arrow[d, hook, "i"'] & X \arrow[r, "f"] \arrow[d, hook] & Y \arrow[d, hook] \\
S'_0 \arrow[r] \arrow[d, "\in W"'] & X'_0 \arrow[r, "f'_0", "W"'] \arrow[d, "W"] & Y'_0 \arrow[d, "W"] \\
S' \arrow[r] & X' \arrow[r, "f'"] & Y'.
\end{tikzcd}
\]
Comme $i$ est une $W$-cofibration et $f \in W$, on a $f'_0 \in W$. D'autre part, comme $S \longrightarrow X$, donc $S'_0 \longrightarrow X'_0$ est
\printedpage{148}
% head: A. Grothendieck : Dérivateurs, Ch. XII, Édition des Universités de Montpellier II et Paris VII
% foot: Version 30 mars 2003
\msp{page 232}
une $W$-cofibration, $X'_0 \longrightarrow X'$ est dans $W$, et pour la même raison, $Y'_0 \longrightarrow Y$ l'est aussi. Donc $f'$ l'est aussi. % [sic] Y'_0 \longrightarrow Y -> Y'_0 \longrightarrow Y', the diagram has Y'_0 \to Y'

\textbf{NB} Ici, c'est un argument formel, valable (sous forme duale) dans tout $(\mathcal{M}, W)$ où on a un énoncé de factorisation pour une flèche quelconque $f : X \longrightarrow Y$ en $X \xrightarrow{i_f} \bar{X} \xrightarrow{\alpha} Y$, avec $i_f \in W$ et $\alpha$ une $W$-fibration. Alors, si $X$, $Y$ sont $W$-fibrés sur $S$, toute $S$-flèche de $X$ dans $Y$ qui est dans $W$, est dans $W^{\mathrm{u}}_S$. Et dualement, si toute flèche $X \longrightarrow Y$ se factorise en $X \longrightarrow \tilde{Y} \longrightarrow Y$, avec $X \longrightarrow \tilde{Y}$ une $W$-cofibration et $\tilde{Y} \longrightarrow Y$ dans $W$.

Il me vient ici une (petite) inquiétude : ai-je au moins prouvé que si
\[
S \longrightarrow X
\]
est une $W$-cofibration, alors pour toute flèche $S' \longrightarrow S$ dans Cat, le carré cocartésien dans Cat % [sic] S' \longrightarrow S -> S \longrightarrow S', the square below has S \to S'
\[
\begin{tikzcd}
S \arrow[r] \arrow[d] & X \arrow[d] \\
S' \arrow[r] & X'
\end{tikzcd}
\]
est aussi $W$-cocartésien (\footnote[126]{dans le sens ancien, spécial à Cat (déf. 3, p. 131). Pour le comparer avec la théorie générale, cf. prop. 20, p. 262.})? Autrement,
\msp{page 233}
à quoi ce carré-là serait-il bon?! Pour le prouver, factorisons encore $S \longrightarrow S'$ en
\[
\underset{\substack{\text{immersion}\\ \text{ouverte de}\\ \text{Dwyer}}}{S \hookrightarrow} S'_0 \underset{\in W}{\longrightarrow} S',
\]
d'où le diagramme de 2 carrés cocartésiens
\[
\begin{tikzcd}
S \arrow[r, "\mathrm{Cof}_W"] \arrow[d, hook, "D_W"'] & X \arrow[d, hook, "D_W"] \\
S'_0 \arrow[r, "\mathrm{Cof}_W"] \arrow[d, "\in W"'] & X'_0 \arrow[d, "\in W"] \\
S' \arrow[r, "\mathrm{Cof}_W"] & X',
\end{tikzcd}
\]
où $X'_0 \longrightarrow X'$ est $\in W$ d'après l'hypothèse sur $X$ sous $S$. Il faut donc vérifier :
\printedpage{149}
% head: A. Grothendieck : Dérivateurs, Ch. XII, Édition des Universités de Montpellier II et Paris VII
% foot: Version 30 mars 2003
\begin{enumerate}
\item[a)] Un carré cocartésien construit sur une flèche qui est \texttt{[une]} immersion de Dwyer, est $W$-cocartésien.
\item[b)] Un carré commutatif dans lequel deux côtés opposés sont dans $W$, est $W$-cocartésien.
\item[c)] Le composé de deux carrés $W$-cocartésiens est $W$-cocartésien.
\end{enumerate}

Admettons provisoirement b) et c), et prouvons a), pour un carré cocartésien
\[
\begin{tikzcd}
U \arrow[r, hook, "i"] \arrow[d, "p"'] & X \arrow[d, "q"] \\
U' \arrow[r, hook, "i'"] & X',
\end{tikzcd}
\]
où $i$ est immersion ouverte de Dwyer. On procède encore comme tantôt, en factorisant $p$ en $U \overset{j}{\hookrightarrow} U'_0 \xrightarrow{p'} U'$, où $j$ est \texttt{[une]} immersion ouverte de Dwyer, $p' \in W$. On a donc un diagramme de deux carrés cocartésiens accolés
\msp{page 234}
\[
\begin{tikzcd}
U \arrow[r, hook, "i \in D_W"] \arrow[d, hook, "D_W"'] & X \arrow[d, hook] \\
U'_0 \arrow[r, hook, "D_W"] \arrow[d, "p' \in W"'] & X'_0 \arrow[d, "q' \in W"] \\
U' \arrow[r, hook] & X',
\end{tikzcd}
\]
où le premier carré est formé d'immersions ouverts. On sait alors qu'il est $W$-cocartésien. Dans le deuxième, on a $p' \in W$, donc $q' \in W$ puisque $i$ est une immersion de Dwyer, donc une $W$-cofibration. Donc il est également $W$-cocartésien. On conclut alors grâce à c) ci-dessus. % [sic] immersions ouverts -> immersions ouvertes, agreement

\noindent Prouvons b), pour
\begin{tikzcd}[baseline=(\tikzcdmatrixname-1-1.base)]
X_0 \arrow[r, "i_1"] \arrow[d, "i_2 \in W"'] & X_1, \arrow[d, "f_1 \in W"] \\
X_2 \arrow[r, "f_2"] & Y
\end{tikzcd}
en l'insérant dans
\begin{tikzcd}[baseline=(\tikzcdmatrixname-1-1.base)]
X_0 \arrow[r, "i_1"] \arrow[d, "i_2 \in W"'] & X_1, \arrow[d, "j_1"] \arrow[ddr, "f_1"] & \\
X_2 \arrow[r, "j_2"] \arrow[drr, "f_2"'] & \mathcal{X} \arrow[dr, "\varphi"] & \\
& & Y
\end{tikzcd}
où $\mathcal{X} = \displaystyle\int X_0$
\begin{tikzcd}[baseline=(\tikzcdmatrixname-2-1.base), row sep=small]
& X_1 \\
{} \arrow[ur] \arrow[dr] & \\
& X_2
\end{tikzcd}
.

Comme $i_2 \in W$, on sait qu'il en résulte $j_1 \in W$, et comme $f_1 \in W$, il en résulte que $\varphi \in W$, q.e.d.
\end{document}
